%! TEX program = xelatex
\documentclass[a4paper,11pt]{report}

\usepackage{./../packages/mainstyle}
\usepackage{./../packages/colors}
\usepackage{./../packages/frameboxes}
\usepackage{./../packages/title}
\usepackage{./../packages/packs}
\usepackage{./../packages/macros}

\usepackage{tabularx} 
\usepackage{array} 
\usepackage[table]{xcolor}

\usepackage[italian]{babel}

\usepackage[T1]{fontenc}
\usepackage{FiraMono} 

\usepackage{listings}
\usepackage{xcolor}

\tcbuselibrary{skins, listings}
\definecolor{hsbg}{HTML}{F8F9FA}       
\definecolor{hspurple}{HTML}{5E5086}   
\definecolor{hsgreen}{HTML}{118C4F}   
\definecolor{hsred}{HTML}{C0392B}      
\definecolor{hsnumber}{HTML}{ADB5BD}  
\definecolor{boxtitlebg}{HTML}{2C3E50} 

\lstdefinestyle{haskellstyle}{
    language=Haskell,
    commentstyle=\color{hsgreen}\itshape,
    keywordstyle=\color{hspurple}\bfseries,
    numberstyle=\tiny\color{hsnumber},
    stringstyle=\color{hsred},
    basicstyle=\ttfamily\footnotesize, 
    breaklines=true,
    showstringspaces=false,
    tabsize=2,                        
    numbers=none                     
}

\newtcblisting{haskellbox}[1]{
    enhanced,
    skin=bicolor,
    colback=hsbg,
    colbacktitle=boxtitlebg,
    coltitle=white,
    colframe=boxtitlebg,
    boxrule=0.5pt,
    arc=4pt,
    drop shadow=black!10!white,
    fonttitle=\bfseries\sffamily,
    title=#1,
    listing only,
    listing options={style=haskellstyle, numbers=left, numbersep=10pt}, % Add numbers here!
    left=15pt,
    top=4pt,
    bottom=4pt,
    breakable
}

% for inline code
\definecolor{inlinebg}{HTML}{F3F4F6}    
\definecolor{inlinetext}{HTML}{24292E}  

% \newcommand{\code}[1]{%
%   \begingroup
%   \setlength{\fboxsep}{2pt}
%   \colorbox{inlinebg}{\ttfamily\color{inlinetext}{\strut #1}}
%   \endgroup
% }

\newtcbox{\code}{
    on line,                           
    fontupper=\ttfamily\color{inlinetext},
    colback=inlinebg,                 
    colframe=inlinebg,               
    boxrule=0pt,                    
    arc=2pt,                       
    boxsep=0pt,                   
    left=3pt,                    
    right=3pt,
    top=1pt,                    
    bottom=1pt
}

\usepackage[autostyle=false]{csquotes}
\MakeOuterQuote{"}

\usepackage{float}
\usepackage{enumitem}
\usepackage{bussproofs}
\usepackage{mdframed}

\usepackage{titlesec}

\setmainfont{XCharter}

\setcoursename{Tecniche di Programmazione\\ Funzionale e Imperativa}
\setcoursebook{slide del professor Salvo}
\setauthorname{aglaia norza}
\setauthoremail{thisisaglaia@gmail.com}
\setauthorgithub{AglaiaNorza}
\setcoverphoto{covermeme.png}

\setcoverwidth{0.6\textwidth}
\setcoverheight{0.6\textwidth}
\setbeforespace{2.5em}
\setafterspace{2.5em}

\newcommand{\apprord}{\sqsubseteq}

\begin{document}

\makecustomcover

\tableofcontents
\newpage

\chapter{Le basi}
\section{Concetti fondamentali}
\begin{itemize}
    \item \textbf{Valutazione lazy:} Haskell valuta le espressioni solo quando è strettamente necessario (permette di gestire agevolmente strutture dati infinite)
    \item \textbf{Trasparenza referenziale:} non essendoci assegnazioni o memoria mutabile (assenza di side-effects), una variabile identifica sempre lo stesso oggetto (permette di manipolare algebricamente i programmi)
    \item \textbf{Polimorfismo e type inference:} molte funzioni operano su tipi generici (es. \code{a}, \code{b}). Il compilatore inferisce automaticamente il \textit{tipo principale} (il più generale possibile); tutti gli altri tipi corretti sono sue istanze.
\item \textbf{Non-terminazione, Bottom e \code{undefined}}: in Haskell, ogni tipo possiede implicitamente un valore speciale (chiamato \textit{bottom}, indicato matematicamente con $\bot$) che rappresenta un'eccezione fatale o una computazione che entra in un loop infinito
    \begin{itemize}
        \item \textbf{la costante \code{undefined}:} è la materializzazione a livello di codice del valore $\bot$ - poiché una computazione che fallisce o non termina può avvenire in qualsiasi contesto (mentre si aspetta un numero, un booleano, ecc.), \code{undefined} possiede il tipo più generale e polimorfo possibile: \code{a} (ovvero $\forall \alpha . \alpha$)
        \item \textbf{Il limite del lambda calcolo (\textit{Occurs check}):} nel lambda calcolo puro, il loop infinito per eccellenza è il combinatore $\Omega = (\lambda x.xx)(\lambda x.xx)$. In Haskell, l'auto-applicazione $\lambda x.xx$ \textbf{non è tipabile}. Perché $x$ possa prendere se stesso come argomento, il suo tipo $\alpha$ dovrebbe essere uguale a una funzione che prende $\alpha$ come parametro e restituisce $\beta$ ($\alpha = \alpha \to \beta$). Questo creerebbe un tipo infinito (durante l'inferenza dei tipi, il compilatore esegue un controllo chiamato \textit{Occurs check} che impedisce alla variabile di tipo $\alpha$ di apparire all'interno della sua stessa definizione, rifiutando di conseguenza l'espressione)
        \item \textbf{Ricorsione esplicita:} poiché il type system ci vieta di creare loop tramite l'auto-applicazione anonima di $\lambda x.xx$, in Haskell la non-terminazione si esprime unicamente attraverso la ricorsione esplicita. Definendo ad esempio \code{omega' = omega'}, creiamo intenzionalmente una computazione divergente che il compilatore accetta, assegnandole il tipo più generale \code{a}.
    \end{itemize}
    \item \textbf{Ideologia funzionale:} l'ideologia funzionale favorisce l'assemblaggio di funzioni semplici tramite composizione (\code{.}). Si pensa in termini di trasformazioni globali di strutture dati, per poi raffinare in un secondo momento i colli di bottiglia prestazionali.
\end{itemize}

\section{Sintassi, Funzioni e Pattern Matching}

\subsection{Definizione e Applicazione}
\begin{itemize}
    \item l'applicazione di funzione: \textbf{associa a sinistra} e si scrive senza virgole (\code{f x y}).
    \item \textbf{aperatori infissi vs prefissi:} le funzioni binarie si usano in modo infisso tra backtick \\ ( \code{10 `div` 2}), gli operatori infissi in modo prefisso tra parentesi ( \code{(+) 3 4}).
\end{itemize}

\subsection{Lambda Espressioni e Combinatori}
Le funzioni anonime si scrivono con backslash (\code{\textbackslash} simula $\lambda$).
\begin{haskellbox}{}
-- identita' (combinatore I)
i :: t -> t
i x = x           -- definizione standard
i' = \x -> x      -- lambda espressione

-- combinatore K (proiettore del primo argomento) (primo assioma di Hilbert)
k :: t1 -> t2 -> t1
k' = \x y -> x

-- combinatore S (sostituzione) (secondo assioma di Hilbert)
s :: (t2 -> t1 -> t) -> (t2 -> t1) -> t2 -> t
s x y z = x z (y z)
\end{haskellbox}

\subsection{Pattern matching}
Il pattern matching decompone le strutture dati in base alla loro forma sintattica. L'ordine delle clausole è strettamente sequenziale (viene scelta la prima dall'alto che fa matching).
\begin{itemize}
    \item \textbf{mancanza di unificazione:} non si possono usare variabili identiche nel pattern per verificarne l'uguaglianza (es. \code{f x x = ...} è invalido)
    \item \textbf{underscore (\code{\_}):} variabile anonima ("don't care") per valori non rilevanti
    \item \textbf{as-pattern (\code{@}):} permette di dare un nome all'intera struttura mentre la si decompone (es. \code{xs@(x:txs)})
\item \textbf{guardie (\code{|}):} permettono di definire il comportamento di una funzione valutando sequenzialmente una serie di \textit{predicati booleani} (condizioni logiche); a differenza del pattern matching che controlla la \textit{forma} sintattica dei dati, le guardie valutano \textit{valori e relazioni} (es. \code{x > 0}, \code{x == y})
    \begin{itemize}
    \item l'ultima guardia è tipicamente \code{otherwise} (che nel Prelude di Haskell è semplicemente un alias per \code{True}), che funge da caso di default ("catch-all") per assicurarsi che la funzione restituisca sempre un valore.
    \end{itemize}
\begin{haskellbox}{}
confronta x y
  | x == y    = "i numeri sono uguali" -- supera la mancanza di unificazione
  | x > y     = "il primo e' maggiore"
  | otherwise = "il secondo e' maggiore"
\end{haskellbox}{}
\item \textbf{equazioni multiple (definizione per casi):} il modo più comune di fare pattern matching è definire la funzione scrivendo più "intestazioni" separate; il compilatore le prova dall'alto verso il basso: se gli argomenti combaciano con la forma (pattern) della riga, esegue quell'equazione, altrimenti controlla la successiva.
\begin{haskellbox}{}
    myAnd True True = True
    myAnd   _   _   = False
\end{haskellbox}{}
\begin{itemize}
    \item è meno comune, ma è possibile fare pattern-matching anche nelle \textit{lambda-espressioni}:
        \begin{haskellbox}{}
            myFst' = \(x, y) -> x
        \end{haskellbox}
\end{itemize}
\end{itemize}

\subsection{Scope locale: \code{let} e \code{where}}
\begin{itemize}
    \item \textbf{\code{let ... in}:} costrutto standalone per definire variabili/funzioni locali (utile per strutturare/destrutturare tuple nelle ricorsioni)
    \item \textbf{\code{where}:} clausola che qualifica la parte destra di una definizione appena conclusa
\end{itemize}
\begin{haskellbox}{}
fibAux n = let (f, fPrec) = fibAux (n-1) in (f + fPrec, f)
-- equivalente a:
fibAux' n = (f + fPrec, f) where (f, fPrec) = fibAux' (n-1)
\end{haskellbox}

\begin{gframe}{differenza tra \code{let} e \code{where}}
   una clausola \code{where} è visibile da tutte le guardie di una specifica equazione. Un \code{let} definito dopo il simbolo \code{=} di una guardia è visibile solo all'interno di quella singola guardia. Per condividere un calcolo tra più guardie, il \code{where} è obbligatorio.
\end{gframe}

\begin{haskellbox}{}
-- condivisione tra guardie
analizzaRaggio r
  | area > 100 = "cerchio grande"
  | area < 10  = "cerchio piccolo"
  | otherwise  = "cerchio medio"
  where area = pi * r^2  -- 'area' calcolata una volta e visibile a tutti i pipe

-- inline: (serve il let perche' non c'e' una dichiarazione di funzione)
volumeCilindro r h = (let area = pi * r^2 in area) * h
\end{haskellbox}

\newpage
\section{Sistema dei Tipi e Classi}

Haskell è \textbf{statically type-safe}: se un programma è tipato correttamente, nessun errore di tipo si verifica durante l’esecuzione.

\subsection{Type Signatures}
In Haskell, grazie al meccanismo di type inference, il compilatore è in grado di dedurre automaticamente il tipo più generale per quasi ogni espressione. Tuttavia, è considerata un'ottima pratica (e in alcuni casi più complessi è strettamente necessario) dichiarare esplicitamente il tipo di una funzione subito sopra la sua definizione. 

Questo si fa usando l'operatore \code{::} (che si legge "ha tipo"). Dichiarare esplicitamente i tipi è utile ai fini della leggibilità del codice e permette di "forzare" il tipo desiderato, aiutando il compilatore a individuare eventuali errori logici in modo molto più preciso.

\begin{haskellbox}{}
-- dichiarazione esplicita del tipo: prende due Int e restituisce un Int
somma :: Int -> Int -> Int
somma x y = x + y
\end{haskellbox}

\subsection{Type Classes}
Le classi limitano il polimorfismo imponendo vincoli (o "interfacce") sui tipi. Più che semplici insiemi, rappresentano classificazioni di tipo algebrico: raggruppano i tipi su cui sono definite certe operazioni. Le classi possono formare gerarchie, dove una sottoclasse "estende" la superclasse ereditandone le operazioni e aggiungendone di nuove.

\begin{itemize}
    \item \textbf{Eq (Equality Types):} comprende i tipi che ammettono l'uguaglianza (\code{==}) e la sua negazione (\code{/=}). Tutti i tipi base (Bool, Int, Float, ecc.) sono istanze di Eq, così come le liste e tuple costruite su tipi Eq
        \subitem \textit{definizione interna di Eq}
        \begin{haskellbox}{}
class Eq a where
  (==), (/=) :: a -> a -> Bool
  -- le classi possono contenere codice (metodi di default)
  x /= y = not (x == y)
\end{haskellbox}

    \item \textbf{Ord (Ordered Types):} sottoclasse di \code{Eq} - tipi ordinabili che supportano operazioni come \code{<=}, \code{<}, \code{>=}, \code{>}, \code{min} e \code{max}.
\begin{itemize}
    \item per definire un'istanza di Ord, è sufficiente fornire il codice per l'operatore \code{<}, il resto può essere interdefinito.
\end{itemize}

    \item \textbf{Num:} la classe più generica per i tipi numerici. È sottoclasse di \code{Eq} e \code{Show}. Implementa operazioni aritmetiche base (\code{+}, \code{-}, \code{*}), oltre a valore assoluto (\code{abs}) e segno (\code{signum}).
        \subitem i numeri interi costanti sono polimorfi rispetto alla classe \code{Num}:
        \begin{haskellbox}{}
> :t 3
3 :: Num a => a  -- 3 puo' essere un Int, Integer, Float, ma non un Char
\end{haskellbox}

    \item \textbf{Altre classi utili:}
        \begin{itemize}
            \item \textbf{Show:} tipi che possono essere stampati (convertiti in \code{String}) tramite il metodo \code{show}. Se provate a stampare a video il risultato di una funzione che non appartiene a \code{Show} (es. una lambda anonima come \code{\textbackslash x -> x}), l'interprete darà errore.
            \item \textbf{Read:} permette di trasformare una \code{String} in un valore tramite il metodo \code{read}
            \item \textbf{Integral:} sottoclasse di \code{Num} per i numeri interi (es. Int, Integer); aggiunge le operazioni di divisione intera \code{div} e modulo \code{mod}
            \item \textbf{Fractional:} sottoclasse di \code{Num} per i numeri frazionari (es. Float, Rational); aggiunge l'operatore di divisione reale \code{(/)} e \code{recip}
        \end{itemize}
\end{itemize}

\subsubsection{Vincoli di classe}
I vincoli di classe si indicano prima del tipo vero e proprio, separati dalla doppia freccia \code{=>}. Vincoli multipli si racchiudono tra parentesi e indicano che la funzione opera su un tipo che deve soddisfare tutte le classi elencate.

\begin{haskellbox}{}
-- esempio: mcd richiede che il tipo 'a' supporti sia l'ordine (<, ==) 
-- sia le operazioni aritmetiche (-)
mcd :: (Ord a, Num a) => a -> a -> a
mcd x y
  | x == y    = x
  | x < y     = mcd (y - x) x
  | otherwise = mcd (x - y) y
\end{haskellbox}

\subsection{Definizione di Tipi: \code{type} e \code{data}}
\begin{itemize}
    \item \textbf{\code{type} (sinonimi):} danno nuovi nomi a tipi esistenti (es. \code{type Casella = (Int, Int)})   
    \item \textbf{\code{data} (tipi algebrici):} costruiscono nuovi tipi tramite costruttori finiti, parametrici o ricorsivi. Possono essere dichiarati istanze di classi.
\end{itemize}

\begin{haskellbox}{}
data SetteNani = Eolo | Pisolo | Brontolo          -- finiti
data Either a b = Left a | Right b                 -- unioni disgiunte
data Maybe a = Just a | Nothing                    -- gestione fallimenti
data List a = Nil | Cons a (List a)                -- ricorsivi/induttivi
data BTree a = Foglia a | Radice (BTree a) (BTree a)
\end{haskellbox}
\subsubsection{Il costruttore di tipo \code{Maybe}}
Il tipo \code{Maybe} è il tipo di computazioni che possono fallire. In Haskell, è il tipo con cui si rappresentano comunemente le eccezioni in maniera sicura e controllata dal sistema di tipi (corrisponde a \code{optional})

Un valore di tipo \code{Maybe} può assumere due forme:
\begin{itemize}
    \item \code{Just v}: contiene un valore \code{v} (la computazione ha avuto successo e restituisce qualcosa)
    \item \code{Nothing}: non contiene nulla (rappresenta una computazione indefinita o fallita)
\end{itemize}

Si può utilizzare per le \textit{funzioni parziali} (che non sono definite per tutti i possibili valori di input (eg divisione per zero)). Tramite il pattern matching, possiamo "spacchettare" il dato per estrarne il contenuto.

\begin{haskellbox}{}
-- definizione del tipo predefinito
data Maybe a = Just a | Nothing

-- 1. costruzione di un Maybe: divisione sicura
-- Se il divisore e' 0 restituisce Nothing, altrimenti restituisce Just il risultato
safediv :: Integral a => a -> a -> Maybe a
safediv n 0 = Nothing
safediv n m = Just (n `div` m)

-- 2. destrutturare un Maybe: funzione "inversa" per estrarre il valore
extract :: Maybe a -> a
extract (Just n) = n
extract Nothing  = undefined 

-- extract (safediv 30 0) -> *** Exception: Prelude.undefined 
\end{haskellbox}

\subsection{Istanze e Derivazione (\code{deriving})}
La keyword \code{deriving} è un meccanismo automatico che permette al compilatore di generare le istanze per le classi predefinite senza dover scrivere codice manuale. 

\begin{itemize}
    \item \textbf{Eq:} genera un'uguaglianza basata sulla sintassi
    \item \textbf{Ord:} deriva l'ordine lessicografico in base alla sequenza dei costruttori
    \item \textbf{Show/Read:} la conversione usa i nomi dei costruttori come stringhe
    \item \textbf{Enum:} permette di generare liste per enumerazione (es. \code{[Lun..Ven]})
\end{itemize}

\begin{haskellbox}{}
-- istanze automatiche tramite deriving 
data Giorni = Lun | Mar | Mer | Gio | Ven | Sab | Dom
    deriving (Eq, Ord, Show, Enum, Read)

-- esempio di istanza manuale 
instance Eq Bool where
    False == False = True
    True == True   = True
    _ == _         = False
\end{haskellbox}

\subsection{Tipabilità e Let-Polymorphism}
\begin{itemize}
    \item Non tutto è tipabile in Haskell. Il classico esempio è l'auto-applicazione pura:
    \begin{itemize}
        \item \code{omega = \textbackslash x -> x x} non è tipabile perché richiederebbe un tipo infinito ($t = t \to t_1$), portando a una teoria dei tipi non decidibile (e quindi a tipi non inferibili dal compilatore)
    \end{itemize}
    
    \item Dal punto di vista del calcolo, \code{let x = N in M} è equivalente all'applicazione di una lambda: \code{(\textbackslash x -> M) N}. Tuttavia, il compilatore Haskell tratta i due casi in modo diverso durante l'inferenza dei tipi.

    \item \textbf{Let-Polymorphism:}

    Il costrutto \code{let} permette di assegnare a una variabile un tipo polimorfo che viene "generalizzato" prima di essere usato nel corpo (\code{in}).
    \begin{itemize}
        \item \textbf{tipabile:} \code{let x = \textbackslash y -> y in x x} è tipabile.

            Haskell vede che \code{x} è l'identità ($a \to a$) e, nel corpo \code{x x}, può istanziare il primo \code{x} come $(a \to a) \to (a \to a)$ e il secondo come $(a \to a)$.
        \item \textbf{non tipabile:} \code{(\textbackslash x -> x x) (\textbackslash y -> y)} \textbf{non} è tipabile. 

            Qui il compilatore deve tipare la funzione \code{\textbackslash x -> x x} prima di sapere che \code{x} sarà l'identità. Come visto sopra, \code{\textbackslash x -> x x} è intrinsecamente non tipabile.
    \end{itemize}

\end{itemize}

\begin{lemmaframe}{Terminazione}{}
    Tutti i lambda-termini tipabili nel sistema dei tipi di Haskell (che non usino ricorsione esplicita tramite equazioni) \textbf{terminano}.
\end{lemmaframe}
 

\subsection{Semantica Denotazionale: ricorsione e punti fissi}

Haskell è un linguaggio dichiarativo: una definizione non è un comando, ma un'equazione. 
Quando definiamo una funzione ricorsiva $f$, stiamo scrivendo un'equazione del tipo $f = \mathcal{E}(f)$, dove $\mathcal{E}$ è un'espressione che contiene $f$ stessa.

\begin{itemize}
    \item \textbf{Bottom ($\bot$)} rappresenta il "non-valore". Indica un calcolo che non termina (loop infinito) o che fallisce (errore). In Haskell, ogni tipo $t$ contiene implicitamente $\bot$.
    \item Se la definizione è un'equazione, la "semantica" (ovvero la funzione reale prodotta dal compilatore) deve essere una soluzione dell'equazione. Matematicamente, una soluzione di $f = T(f)$ è un \textbf{punto fisso} del funzionale $T$.
\end{itemize}

\begin{gframe}{Approfondimento: ricorsione e non-terminazione: \code{omega'}}
Prendiamo la funzione
\begin{haskellbox}{}
omega' x = omega' x
-- > :t omega'
-- omega' :: t -> t1
\end{haskellbox}
Questa funzione ha il tipo più generico possibile (\code{t -> t1}) perché, non terminando mai, non produce mai un valore reale.

Rimuovendo $x$, otteniamo $omega' = I(omega')$, dove $I$ è l'identità.

Qualsiasi funzione $f$ soddisfa $f = I(f)$. Quale scegliere?

\begin{itemize}
    \item Haskell sceglie la funzione che "fa il minimo indispensabile" per soddisfare l'equazione. Si usa l'ordinamento $\apprord$ dove $f \apprord g$ significa che $g$ è "più definita" di $f$.
    \item \textbf{Costruzione di Kleene (least-fixpoint):} 

        il computer "costruisce" la funzione partendo dal nulla ($\bot$):
    \[ f_0 = \bot \quad (\text{non so nulla}) \]
    \[ f_1 = T(f_0) \quad (\text{applico la definizione una volta}) \]
    \[ f_{\infty} = \text{limite della catena} = \textit{lfp}(T) \]
    Nel caso di \code{omega'}, la catena è $\{\bot, \bot, \dots\}$. Il limite è $\bot$. La semantica di \code{omega'} è dunque la funzione che restituisce sempre "non-valore".
\end{itemize}

\end{gframe}

Il punto fisso è quindi la "semantica" perché è l'unico modo per dare un valore univoco e calcolabile a una definizione circolare. Il \textit{minimo} punto fisso garantisce che il significato sia quello "prodotto" effettivamente dal calcolo ricorsivo.

\begin{gframe}{\code{myundefined}}
Mentre \code{omega'} è una \textit{funzione} che non termina, possiamo definire un \textit{valore} che non termina:
\begin{haskellbox}{}
myUndefined = myUndefined
-- > :t myUndefined
-- myUndefined :: a
\end{haskellbox}

\begin{itemize}
    \item  Poiché \code{myUndefined} ha tipo \code{a}, esso \textbf{appartiene a tutti i tipi}. È l'implementazione "home-made" della costante \code{undefined} di Haskell. Rappresenta il valore $\bot$ per ogni possibile tipo del linguaggio.
    \item Matematicamente, \code{myUndefined} è il \textbf{punto fisso dell'identità}: 
\begin{enumerate}
\item la definizione \code{x = x} corrisponde all'equazione funzionale $x = I(x)$
\item la soluzione di questa equazione è, per definizione, il punto fisso dell'identità $I$
\item poiché Haskell usa la semantica del minimo punto fisso, e l'identità applicata al valore nullo ($\bot$) restituisce ancora $\bot$, il risultato è la non-terminazione
\item questo spiega perché \code{myUndefined} ha tipo \code{a}: non essendo "sceso" verso nessun valore specifico (come un intero o una stringa), rimane una pura astrazione di errore/loop che può "fingere" di essere qualsiasi tipo.
\end{enumerate}
\end{itemize}



Haskell permette quindi di manipolare la non-terminazione come un valore reale. Ogni tipo in Haskell è in realtà "puntato" (\textit{pointed}), ovvero contiene sempre almeno un valore: $\bot$.
\end{gframe}




\section{Strutture Dati Predefinite}

\subsection{Tipi di Base e Tuple}
\begin{itemize}
    \item \textbf{Booleani:} \code{Bool} (\code{True}, \code{False}). Operatori: \code{not}, \code{\&\&}, \code{||}.
    \item \textbf{Interi:} \code{Int} (dimensione fissa), \code{Integer} (precisione illimitata).
    \item \textbf{Tuple:} n-uple fisse e disomogenee (es. \code{(String, Bool)})
        \subitem offrono le funzioni \textit{su coppie}: \code{fst}, \code{snd} (prendono primo e secondo argomento)
\end{itemize}

\subsection{Liste}
Sequenze di elementi omogenei. Definite dal costruttore vuoto \code{[]} e dall'operatore infisso \textit{cons} \code{(:)}. Le stringhe sono \code{[Char]}. 

Zucchero sintattico: \code{[1,2,3]} equivale a \code{1:2:3:[]}.

\begin{haskellbox}{}
-- enumerazioni 
[m..n]           -- range finito (lista con elementi da m a n)
[m, n..p]        -- range con step (es. [1, 3..11] dispari)
[m..]            -- lista infinita (lazy)

\end{haskellbox}
Haskell supporta le \textbf{list comprehension} (definizioni set-like del contenuto di una lista)

\begin{haskellbox}{}
-- list comprehension (spesso tradotte internamente in map/filter)
[x^2 | x <- [1..5]]               -- quadrati
[x | x <- xs, x `mod` 2 == 0]     -- filtri / guardie
[(x,y) | x <- xs, y <- ys]        -- prodotto cartesiano
\end{haskellbox}

\subsubsection{Funzioni su Liste}
\begin{itemize}
    \item \textbf{Estrazione:} \code{head}, \code{tail}. \code{last}, \code{init}. \code{(!!)} per estrarre all'indice $n$
        \begin{haskellbox}{}
head [1, 2, 3]  -- risultato: 1
tail [1, 2, 3]  -- risultato: [2, 3]
last [1, 2, 3]  -- risultato: 3
init [1, 2, 3]  -- risultato: [1, 2]
[1, 2, 3] !! 1  -- risultato: 2
\end{haskellbox}{}

    \item \textbf{Ispezione:} \code{null}, \code{length}
\begin{haskellbox}{}
null []          -- risultato: True
length [1, 2, 3] -- risultato: 3
\end{haskellbox}{}

    \item \textbf{Taglio:} \code{take n}, \code{drop n}, \code{splitAt n} (tupla con take e drop).
\begin{haskellbox}{}
take 2 [1, 2, 3, 4]    -- risultato: [1, 2]
drop 2 [1, 2, 3, 4]    -- risultato: [3, 4]
splitAt 2 [1, 2, 3, 4] -- risultato: ([1, 2], [3, 4])
\end{haskellbox}{}

    \item \textbf{Generali:} \code{++} (concatenazione), \code{reverse}
\begin{haskellbox}{}
[1, 2] ++ [3, 4]  -- risultato: [1, 2, 3, 4]
reverse [1, 2, 3] -- risultato: [3, 2, 1]
\end{haskellbox}{}

    \item \textbf{Matematica/logica:} \code{sum}, \code{minimum}, \code{maximum}. \code{and}, \code{or}. \code{any p}, \code{all p} (predicati su almeno uno o tutti)
        \begin{haskellbox}{}
sum [1, 2, 3]       -- risultato: 6
minimum [5, 2, 9]   -- risultato: 2
maximum [5, 2, 9]   -- risultato: 9
and [True, False]   -- risultato: False
or [True, False]    -- risultato: True
any (> 2) [1, 2, 3] -- risultato: True (almeno uno è > 2)
all (> 0) [1, 2, 3] -- risultato: True (tutti sono > 0)
\end{haskellbox}{}

    \item \textbf{Liste annidate/multiple:} \code{concat} (appiattisce liste di liste), \code{zip} (unisce liste in coppie, fermandosi alla più corta)
\begin{haskellbox}{}
concat [[1, 2], [3, 4]]    -- risultato: [1, 2, 3, 4]
zip [1, 2] ['a', 'b', 'c'] -- risultato: [(1, 'a'), (2, 'b')]
\end{haskellbox}{}
\end{itemize}

\section{Funzion(al)i di Ordine Superiore}
I funzionali sono funzioni che prendono altre funzioni come argomento o le restituiscono come risultato. Favoriscono la \textbf{composizionalità}, permettendo di costruire programmi complessi tramite piccoli blocchi riutilizzabili.

\begin{itemize}
    \item \textbf{curryficazione:} si fonda sull'isomorfismo $A \times B \rightarrow C \cong A \rightarrow (B \rightarrow C)$. Le versioni curryficate permettono di passare un solo argomento alla volta, favorendo l'applicazione parziale.
    \item \textbf{$\eta$-regola:} è possibile omettere un parametro in una definizione se esso compare identico alla fine di entrambi i lati dell'equazione.
\end{itemize}

\begin{haskellbox}{}
-- composizione (.) : Equivale a f(g(x))
(.) f g x = f (g x)

-- curry e uncurry (Isomorfismo A x B -> C  < = = >  A -> (B -> C))
curry :: ((a, b) -> c) -> a -> b -> c
uncurry :: (a -> b -> c) -> (a, b) -> c

-- flip: inverte l'ordine dei parametri
flip :: (a -> b -> c) -> b -> a -> c
myFlip f a b = f b a

-- flip (:) [2, 3] 1
-- [1, 2, 3]

-- map e filter
map :: (a -> b) -> [a] -> [b]         -- proprieta': map (f . g) = map f . map g
filter :: (a -> Bool) -> [a] -> [a]

-- zipWith: applica f. binaria agli elementi di due liste
zipWith :: (a -> b -> c) -> [a] -> [b] -> [c]
\end{haskellbox}

\subsection{Famiglia dei \code{fold} (schemi di ricorsione)}
I \code{fold} astraggono il concetto di ricorsione strutturale sulle liste, sostituendo i costruttori della struttura dati (l'operatore \code{(:)} e la lista vuota \code{[]}) con funzioni e valori specifici. Se l'operatore \code{\#} è associativo e ha \code{e} come elemento neutro, allora $\code{foldl (\#) e} = \code{foldr (\#) e}$.

\begin{itemize}
    \item \textbf{\code{foldr} (right fold):} associa a destra; la riduzione avviene "al ritorno" dalle chiamate ricorsive (costruisce espressioni). 
\begin{itemize}
    \item sostituisce ricorsivamente il costruttore \code{(:)} con la funzione \code{f} data e \code{[]} con il valore base \code{z}.
        \item espansione: \code{foldr f z [x1, x2, x3]} diventa \code{x1 `f` (x2 `f` (x3 `f` z))}
        \item può terminare anche su liste infinite se la funzione \code{f} non è stretta nel suo secondo argomento (ovvero se non ha sempre bisogno di valutare la ricorsione per restituire un risultato).
\end{itemize}
        \begin{haskellbox}{}
-- somma con foldr
foldr (+) 0 [1, 2, 3]
-- valutazione: 1 + (2 + (3 + 0)) = 6

-- implementazione di map usando foldr
myMap f xs = foldr (\x acc -> f x : acc) [] xs

\end{haskellbox}

    \item \textbf{\code{foldl} (left fold):} associa a sinistra; simula un ciclo iterativo passando un accumulatore "in avanti"
\begin{itemize}
    \item valuta l'espressione partendo da sinistra e procedendo verso l'interno - valuta tutto solo alla fine della lista
        \item espansione: \code{foldl f z [x1, x2, x3]} diventa \code{((z `f` x1) `f` x2) `f` x3}
        \item \textit{nota bene:} In \code{foldl}, la funzione prende come primo argomento l'accumulatore, e come secondo l'elemento corrente della lista.
\end{itemize}
        \begin{haskellbox}{}
-- sottrazione con foldl
foldl (-) 0 [1, 2, 3]
-- valutazione: ((0 - 1) - 2) - 3 = -6

-- sottrazione con foldr per confronto
foldr (-) 0 [1, 2, 3]
-- valutazione: 1 - (2 - (3 - 0)) = 2

-- reverse usando foldl
myReverse xs = foldl (\acc x -> x : acc) [] xs
\end{haskellbox}

    \item \textbf{\code{foldr1} e \code{foldl1}:} varianti che non richiedono un valore iniziale esplicito, poiché usano il primo o l'ultimo elemento come base; \textit{danno errore su lista vuota}.
\begin{haskellbox}{}
-- utile quando non esiste un elemento neutro sensato (es. min/max assoluto)
myMinimum = foldr1 min
myMinimum [5, 3, 8] -- risultato: 3
-- myMinimum []     -- genera errore a runtime
\end{haskellbox}{}
\end{itemize}

\section{Alcuni esempi di stile funzionale}
Lo stile funzionale evita accumulatori ed indici, preferendo la manipolazione globale delle strutture dati.

\begin{itemize}
    \item \textbf{prodotto scalare:} può essere implementato in diversi modi (flessibilità dei funzionali):

\begin{haskellbox}{}
-- accoppia gli elementi, trasforma "*" in una funzione che accetta una tupla
-- applica la moltiplicazione agli elementi, e somma i prodotti
ps1 xs ys = sum (map (uncurry (*)) (zip xs ys))

-- crea una lista di funzioni parzialmente applicate
-- (da [5, 4] si ottiene [(5*), (4*)])
-- zApp/applyL prende la lista di funzioni ^ e la applica sulla seconda lista
ps2 xs ys = sum (applyL (map (*) xs) ys)

-- zipWith fa zip e map in un passo solo
ps3 xs ys = sum (zipWith (*) xs ys)
\end{haskellbox}

    \item controllare se una lista è ordinata:
\begin{haskellbox}{}
ordinata xs = and (zipWith (<=) xs (tail xs))
\end{haskellbox}

\item \textbf{ricerca su lista} (\code{findVPH}): si può implementare accoppiando la lista agli indici (\code{zip [0..] xs}) e filtrando per il valore cercato.
\end{itemize}

\chapter{Lambda Calcolo}
Ideato da Alonzo Church nel 1931, il $\lambda$-calcolo nasce per esplicitare il \textit{processo meccanico di calcolo} (sintassi) di una funzione, in contrasto con la visione matematica classica che vede la funzione solo come una relazione statica tra insiemi (semantica).

\section{Sintassi e Computazione}
La sintassi del $\lambda$-calcolo è minimalista ed è composta da tre regole:
\begin{itemize}
    \item \textbf{variabile:} $x, y, z \dots$
    \item \textbf{astrazione:} $\lambda x. M$ (corrisponde alla definizione di una funzione con parametro $x$ e corpo $M$)
    \item \textbf{applicazione:} $(M N)$ (corrisponde all'applicazione della funzione $M$ all'argomento $N$)
\end{itemize}
\vspace{0.5em}
in BNF, quindi, la grammatica del lambda-calcolo è quindi:
$$T ::= V \mid \lambda V.T \mid (T \ T)$$

\begin{gframe}{Regole di associazione}
    \begin{itemize}
        \item \textbf{l'applicazione associa a sinistra:} $F N_1 N_2 \dots N_n$ equivale a $(\dots (F N_1) \dots N_n)$
        \item \textbf{l'astrazione associa a destra:} $\lambda x_1 \dots x_n. M$ equivale a $\lambda x_1. (\lambda x_2. (\dots (\lambda x_n. M)\dots))$
    \end{itemize}
\end{gframe}

La computazione avviene tramite la \textbf{$\beta$-riduzione} (beta-regola): un termine della forma $(\lambda x.M) N$ si dice \textit{redex} e si riduce sostituendo tutte le occorrenze di $x$ in $M$ con il termine $N$.
Formalmente:
$$(\lambda x.M) N \rightarrow_\beta M[N/x]$$

Un termine che non contiene più $\beta$-redex è detto in \textbf{forma normale} e rappresenta un valore finale (il risultato della computazione).

\section{Variabili e Combinatori}
\begin{itemize}
    \item \textbf{Variabili libere ($FV$) e legate:} in $\lambda x. M$, la variabile $x$ è \textit{legata} (bound) all'astrazione. Qualsiasi variabile non legata è \textit{libera} (free).
    \item \textbf{$\alpha$-conversione (alfa-regola):} per evitare la cattura accidentale di variabili libere durante la sostituzione, si possono rinominare le variabili legate (es. $\lambda x.x \equiv \lambda y.y$). 
        \begin{itemize}
            \item è buona pratica seguire la ``convenzione di Barendregt'', che stabilisce che, all'interno di un termine lambda o di un insieme di termini, tutte le variabili legate devono avere nomi distinti tra loro e diversi da tutte le variabili libere presenti nel termine stesso
        \end{itemize} 
    \item  I termini tali che ($FV(M) = \emptyset$) ($\lambda$-termini chiusi), ovvero senza variabili libere, sono chiamati \textbf{combinatori}
\end{itemize}

\subsection{Combinatori noti e strutture di controllo}
Il $\lambda$-calcolo puro non ha tipi di dato predefiniti: tutto (booleani, numeri, if-then-else) deve essere codificato tramite funzioni.

\begin{itemize}
    \item \textbf{Identità ($I$):} $I \equiv \lambda x.x$
    \item \textbf{Cancellatori (booleani):}
        \subitem \textbf{True ($K$):} $K \equiv \lambda x y. x$ (prende due argomenti e restituisce il primo)
        \subitem \textbf{False ($O$):} $O \equiv \lambda x y. y$ (prende due argomenti e restituisce il secondo)
    \item \textbf{If-Then-Else:} grazie alla definizione dei booleani, la struttura condizionale è codificata come un'applicazione: $if\ x\ then\ M\ else\ N \equiv x M N$. Se $x$ è True ($K$) selezionerà $M$, se è False ($O$) selezionerà $N$.
    \item \textbf{Compositori ($S$):} $S \equiv \lambda x y z. x z (y z)$ (si può dimostrare che $S K K \approx I$).
    \item \textbf{Duplicatori e divergenza:} 
        \subitem $\omega \equiv \lambda x.xx$ (auto-applicazione)
        \subitem $\Omega \equiv \omega \omega \rightarrow \omega \omega \rightarrow \dots$ rappresenta la \textbf{funzione indefinita}, il loop infinito (non tipabile)
\end{itemize}

\subsubsection{Numeri di Church e iteratori}
I \textbf{Numerali di Church} rappresentano i numeri naturali come funzioni di ordine superiore che \textit{iterano} un'operazione. Il numero $n$ è una funzione che prende una funzione (es. successore) $s$ e un valore di partenza (es. zero) $z$, e applica $s$ ad $z$ esattamente $n$ volte
$$\underline{n} \equiv \lambda s z. s(s(\dots(sz)\dots)) \quad [n\text{ applicazioni}]$$
\begin{itemize}
    \item $\underline{0} \equiv \lambda s z. z$ (equivale a \code{False})
    \item $\underline{3} \equiv \lambda s z. s(s(s(z)))$
\end{itemize}

Per calcolare il successore ($+1$) di un numero di Church, basta creare una funzione che applica $s$ una volta in più: $succ \equiv \lambda x s z. s(x s z)$ (o anche $\lambda x s z. xs (s z)$)

I numerali di Church non sono quindi solo dati, ma \textbf{iteratori intrinseci}.

\section{Ricorsione e Combinatori di Punto Fisso}
Per definire funzioni ricorsive complesse il $\lambda$-calcolo necessita di un principio di ricorsione generale. 

\begin{thmframe}{Combinatore di Punto Fisso}{}
 Nel $\lambda$-calcolo esiste un termine $Y$ (detto \textbf{Combinatore di Punto Fisso}) tale che, per ogni termine $M$, si ha $Y M \rightarrow M (Y M)$. 
 
 Il combinatore permette di trovare quindi il punto fisso di una funzione (quel valore $x$ tale che $f(x) = x$)
\end{thmframe}

Due combinatori di punto fisso sono:
\begin{itemize}
    \item \textbf{Combinatore di Turing ($Y_T$):} definito come $Y_T = \theta \theta$ dove $\theta \equiv \lambda x y. y(x x y)$.
\begin{align*}
Y_T M &\equiv (\theta \theta) M \\
&\equiv ((\lambda x y. y(x x y)) \theta) M \\
&\rightarrow_\beta (\lambda y. y(\theta \theta y)) M \\
&\rightarrow_\beta M (\theta \theta M) \\
&\equiv M (Y_T M)
\end{align*}

    \item \textbf{Combinatore di Curry ($Y_C$):} definito come $Y_C \equiv \lambda f. (\lambda x.f(xx))(\lambda x.f(xx))$. Se applicato all'identità $I$, questo combinatore riduce a $\Omega$ (divergenza).
    \begin{align*}
    Y_C M &= (\lambda f. (\lambda x. f(x x)) (\lambda x. f(x x))) M \\
    &\rightarrow_\beta (\lambda x. M(x x)) (\lambda x. M(x x)) \\
    &\rightarrow_\beta M ( (\lambda x. M(x x)) (\lambda x. M(x x)) ) \\
    &\equiv M (Y_C M)
    \end{align*}
    {\small \color{gray}(l'ultima equivalenza $\equiv$ vale perché $(\lambda x. M(x x)) (\lambda x. M(x x))$ è la forma ridotta di $Y_C M$)}

\end{itemize}

\begin{thmframe}{Tesi di Church-Turing}{}
 Ogni funzione calcolabile è $\lambda$-definibile.
\end{thmframe}


\chapter{Temi avanzati di programmazione funzionale}

\section{Fusion law per foldr}

La \textbf{Fusion Law} per \code{foldr} è un principio di induzione generalizzato che permette di trasformare la composizione di una funzione con una \code{foldr} in un'unica fold. Essa stabilisce le condizioni per cui vale l'uguaglianza:

\code{f . foldr g a = foldr h b}

L'obiettivo è quindi trovare funzioni o valori \code{h} e \code{b} tali che l'applicazione di \code{f} al risultato di una \code{foldr} sia equivalente ad una nuova \code{foldr} che lavora direttamente su di essi - l'utilità computazionale risiede nella capacità di ``fondere'' due passaggi (una trasformazione \code{f} applicata dopo un fold con \code{g}) in un unico ciclo sulla struttura dati, evitando allocazioni intermedie.

\begin{thmframe}{Fusion law}{}
    Siano \code{f}, \code{g}, \code{h} funzioni e \code{a}, \code{b} valori. L'uguaglianza \code{f . foldr g a = foldr h b} è garantita se sono soddisfatti i seguenti tre requisiti:
    \begin{enumerate}
        \item \code{f} è stretta (\code{f undefined = undefined}) (così che valga anche nel caso di liste parziali/indefinite)
        \item \textbf{caso base}: \code{f a = b} - \code{f} applicata all'accumulatore iniziale della prima \code{foldr} (\code{a}) deve dare come risultato l'accumulatore iniziale della seconda (\code{b})
        \item \textbf{passo induttivo} (condizione di fusione): per ogni \code{x}, \code{y}, deve valere: 
            \subitem \code{f (g x y) = h x (f y)} - in questo modo, l'operazione \code{h} del nuovo \code{foldr} riproduce esattamente l'effetto di \code{g} seguito da \code{f}
    \end{enumerate}

\end{thmframe}

\begin{gframe}{parallelo con un omorfismo}
    La fusion law può essere interpretata come la dimostrazione che \code{f} agisce come un \textbf{omomorfismo} tra due strutture algebriche di calcolo:
\begin{itemize}
    \item \textbf{preservazione dell'identità}: \code{f a = b} mappa l'elemento neutro (accumulatore iniziale) del primo dominio in quello del secondo, analogamente a $f(1_A) = 1_B$.
    \item \textbf{preservazione dell'operazione}: La condizione \code{f(g x y) = h x (f y)} ricalca la proprietà omomorfica $f(a \bullet b) = f(a) \circ f(b)$. 
    \begin{itemize}
        \item \textit{sorgente (\code{g x y})}: rappresenta l'operazione di combinazione di un elemento \code{x} con il risultato parziale \code{y} tramite la funzione \code{g}
        \item \textit{destinazione (\code{h x (f y)})}: rappresenta la nuova operazione \code{h} che combina \code{x} con il risultato già trasformato da \code{f}
    \end{itemize}
\item \textbf{NB}: in \code{h x (f y)}, l'elemento \code{x} non viene trasformato da \code{f} - nella struttura della lista, \code{x} è un dato ``guida'' di tipo $A$, mentre \code{y} è il risultato della ricorsione (accumulatore) di tipo $B$. La fusion law mira a trasformare il \textit{risultato complessivo} del fold, non i singoli elementi contenuti nella lista.
\end{itemize}
\end{gframe}

\begin{proofframe}[title=Derivazione formale]
Per dimostrare la validità dell'uguaglianza \code{f . foldr g a = foldr h b}, analizziamo separatamente i due lati della relazione applicando i principi di induzione sulla struttura della lista:

Ricordando la definizione di \code{foldr} basata sulla decomposizione della lista:
\begin{itemize}
    \item \code{foldr g a [] = a}
    \item \code{foldr g a (x:xs) = g x (foldr g a xs)} 
\end{itemize}

\textbf{Caso \code{[]} (base)}:
\begin{itemize}
    \item \textbf{sinistra}: \code{(f . foldr g a) []} \\
        \code{= f (foldr g a [])} \textit{ (def di .)} \\
    \code{= f a} \textit{(def di foldr)}
\item \textbf{destra}: \code{foldr h b []}   \\
          \code{= b} \textit{(def di foldr)}
    \item da cui deriva la prima condizione: \code{f a = b}
\end{itemize}

\textbf{Caso \code{(x:xs)} (induttivo)}:
\begin{itemize}
    \item \textbf{sinistra}: \code{(f . foldr g a) (x:xs)} \\
        \code{= f (foldr g a (x:xs))}  \textit{ (def di .)}\\
        \code{= f (g x (foldr g a xs))} \textit{ (def di foldr)} 
    \item \textbf{destra}: \code{foldr h b (x:xs)} \\
        \code{= h x (foldr h b xs)} \textit{ (def di foldr)} \\
        \code{= h x (f (foldr g a xs))} \textit{ (induzione)}
      \item ponendo \code{y = foldr g a xs}, deriviamo la condizione di fusione:
          \subitem\code{f (g x y) = h x (f y)}
\end{itemize}
\end{proofframe}

\section{TODO: call-by-name + laziness (08)}
\section{TODO: funzioni strette e non strette}
\section{TODO: tecniche di ottimizzazione (09)}

\section{scanl e scanr}
\code{scanl} e \code{scanr} sono due funzionali che permettono di processare una lista restituendo non solo il risultato finale (come fa un fold), ma una lista contenente tutti i \textbf{risultati intermedi} della computazione.

Come \code{foldr} e \code{foldl}, i due funzionali si distinguono per la direzione in cui accumulano i risultati:

\begin{itemize}
    \item \code{scanl} calcola i risultati ``in avanti'', applicando la funzione a tutti i \textbf{prefissi} della lista:
\begin{haskellbox}{}
scanl (#) v [x, y, z] = [v, v#x, (v#x)#y, ((v#x)#y)#z]

scanl (+) 0 [1, 2, 3] = [0, 1 (ovvero 0+1), 3 (0+1+2), 6 (0+1+2+3)]
\end{haskellbox}

\item \code{scanr} calcola i risultati ``all'indietro'', applicando la funzione a tutti i \textbf{suffissi} della lista
\begin{haskellbox}{}
scanr (#) v [x, y, z] = [x#(y#(x#e)), y#(z#e), (z#v), v]

scanr (+) 0 [1, 2, 3] = [6 (3+2+1), 5 (3+2), 3, 0 ]
\end{haskellbox}

\end{itemize}

\begin{gframe}{Derivazione del codice di \code{scanl}}
La specifica formale di \code{scanl} definisce il funzionale come l'applicazione di un \code{foldl} a tutti i segmenti iniziali (\code{inits}) di una lista:

\begin{haskellbox}{}
scanl f e = map (foldl f e) . inits
\end{haskellbox}

Tuttavia, questa definizione non è efficiente (porta ad un comportamento quadratico).
Tramite la program calculation, è possibile derivare una definizione ricorsiva diretta, di complessità lineare ($\Theta(n)$).

Le equazioni che definiscono la versione efficiente (lineare) di \code{scanl} sono le seguenti:
\begin{haskellbox}{}
scanl f e [] = [e]
scanl f e (x:xs) = e : scanl f (f e x) xs
\end{haskellbox}

In questa forma, il risultato dell'applicazione della funzione (\code{f e x}) viene passato come nuovo valore iniziale per la chiamata ricorsiva sulla coda della lista, evitando di ricalcolare i prefissi da zero.
\end{gframe}

\section{Strutture dati infinite}
La valutazione lazy permette di maneggiare in modo astratto ed elegante \textbf{strutture dati infinite}.

Gli abitanti di un tipo ricorsivo in Haskell sono la \textbf{massima algebra chiusa} rispetto all'applicazione dei costruttori.

\begin{defframe}{Massima vs minima algebra chiusa}{}
\begin{itemize}
    \item \textit{Minima algebra chiusa} (induzione): insieme di tutti i dati finiti che si possono costruire partendo dai casi base (esempio: numeri naturali di Peano)
    \item \textit{Massima algebra chiusa} (co-induzione): include anche gli oggetti costruiti applicando i costruttori infinite volte - i tipi di Haskell ammettono questa natura infinita
        \subitem per esempio, la co-algebra dei naturali includerebbe anche il `naturale infinito' 
        $$\omega = S(S(\dots))=S^{\infty}(\dots)$$
\end{itemize}
\end{defframe}

Nel caso delle liste, le liste \textit{finite} sono la minima algebra chiusa rispetto a \code{:} e \code{[]}, mentre le liste infinite (o \textbf{stream}) sono la relativa massima algebra.

Eliminando i costruttori costanti (casi base) dalle algebre induttive, si possono ``isolare'' gli elementi infiniti:
\begin{itemize}
    \item eliminando \code{zero} dai naturali, si ottiene l'algebra che contiene solo il naturale infinito
    \item eliminando \code{[]} dalle liste, si ottiene l'algebra che contiene solo stream
\end{itemize}

\subsection{Definizioni naif e circolari}
Il modo che si sceglie per generare stream ne cambia drasticamente l'efficienza. 

Possiamo definire le implementazioni di stream in due categorie: naif (basate sul ricalcolo), e circolari (basate sulla condivisione)

\subsubsection{Definizioni naif}
Una definizione naif richiama esplicitamente se stessa passando nuovamente gli argomenti

\begin{haskellbox}{esempio: potenze e iterate}
powers n = 1 : map (n*) (powers n)

myIterate f x = x : map f (myIterate f x)
\end{haskellbox}

Queste implementazioni hanno complessità quadratica ($O(n^2)$), in quanto, per calcolare l'$n$-esimo elemento, il programma deve ripercorrere tutti i passaggi precedenti.

La complessità è causata da un'assenza di memoizzazione del lavoro già fatto: ogni volta che si chiede un nuovo elemento viene generata una nuova chiamata alla funzione, creando una catena infinita di funzioni annidate (\code{map (*2) (map (*2) (map (*2) ...))}).


\subsubsection{Definizioni circolari e lineari}
Una definizione si dice circolare quando la struttura dati viene definita in termini di se stessa, creando un puntatore in memoria (Graph Reduction) invece di generare nuove chiamate a funzione. Questo garantisce una complessità lineare ($O(n)$) perché ogni elemento viene calcolato una volta sola.

Esistono due modi principali per implementarle:
\begin{itemize}
    \item \textbf{livello globale}: senza parametri, il nome della lista è sufficiente per creare il riferimento circolare
\begin{haskellbox}{}
ones = 1 : ones
fibs = 0 : 1 : zipWith (+) fibs (tail fibs)
\end{haskellbox}
    \item \textbf{all'interno di funzioni}: si usano \code{where} o \code{let} per ``fissare'' i parametri e definire la lista ricorsivamente rispetto alla variabile locale
\begin{haskellbox}{}
powers' n = ps where ps = 1 : map (*n) ps
myIterate' f x = xs where xs = x : map f xs
\end{haskellbox}
\end{itemize}

È utile saper distinguere tra velocità (linearità) e struttura di memoria (circolarità). Non tutte le ricorsioni dirette sono inefficienti.
\begin{haskellbox}{}
iterate1 f x = x : iterate1 f (f x)
\end{haskellbox}

Questa definizione è lineare ($O(n)$) pur non essendo circolare. Ogni passo produce un elemento applicando la funzione all'argomento corrente, senza accumulare operazioni sospese (come catene di \code{map}) che invece renderebbero la computazione quadratica.

\begin{gframe}{Produttività}
La ``regola d'oro'' per gli stream è garantire la \textbf{produttività}: il processo di generazione deve consumare meno input di quanto output produce. 
\begin{itemize}
    \item esempio non produttivo: \code{incstnts = (head incstnts) : (tail incstnts)}.
    \subitem la definizione richiede di consumare la testa della lista prima ancora di averla generata, portando a una mancata terminazione (il processo ``si mangia la coda'')
\end{itemize}
\end{gframe}

\section{Liste infinite come limiti}
Una lista infinita può essere vista come il \textbf{limite delle sue approssimazioni finite}.

Per gestire l'infinito e il calcolo parziale, Haskell usa il simbolo $\bot$ (bottom, anche usato per la contraddizione in logica), che rappresenta una computazione che non termina o che ha dato errore (corrisponde ad \code{undefined}).

L'esecuzione di un programma può essere vista come un \textbf{accumulo di informazioni}:
\begin{itemize}
    \item un valore non ancora calcolato ($\bot$) non ha informazioni
    \item man mano che il programma avanza, $\bot$ viene sostituito da valori reali (es: $1:\perp$, poi $1:2:\perp$)
\end{itemize}

Possiamo quindi definire una relazione di ordine parziale $\apprord$, l'\textbf{ordine di approssimazione}.

Diciamo che $x \apprord y$ ($x$ approssima $y$) se $y$ contiene almeno tutte le informazioni di $x$.

Per esempio, una lista infinita \code{[1..]} può essere vista come il limite di una sequenza di approssimazioni che diventano sempre più precise:
$$\bot \apprord 1:\bot \apprord 1:2:\bot \apprord 1:2:3:\bot \dots $$

Il limite di questa catena è la lista infinita completa.

(Invece, una lista come $\bot, 1 : \bot, 2 : 1: \bot, 3 : 2 : 1 : \bot, \dots$ non converge a nessun limite: non si stabilizza a nessun valore e offre informazioni contraddittorie).

\subsection{Domini, ordine e limiti}
Vediamo nel dettaglio come Haskell misura l'informazione prodotta da un programma.

\subsubsection{Tipi semplici}
Per quanto riguarda i tipi semplici, l'informazione è un ``tutto o niente'': o il valore è calcolato, o non lo è.

La relazione $\apprord$ per i tipi semplici è \textit{piatta} (\textbf{flat domain}): 
$$x \apprord y \text{ sse } x = \bot \text{ oppure } x = y$$

Si ha quindi per esempio $\bot \apprord$\code{True}, ma non ci sono relazioni tra \code{True} e \code{False}.

\subsubsection{Costruttori di tipo}
L'ordine si estende naturalmente ai tipi composti. L'idea è che una struttura è ``più definita'' di un'altra se i suoi componenti lo sono.

\begin{itemize}
    \item \textit{Coppie}: $(x, y) \apprord (x', y') \iff x \apprord x' \land y \apprord y'$
    \item \textit{Either}: 
        \begin{itemize}
            \item $\bot \apprord$\code{Left}$\bot$, \code{Right}$\bot$
            \item \code{Left}$x \apprord$ \code{Left} $x' \iff x \apprord x'$
            \item \code{Right}$x \apprord$ \code{Right} $x' \iff x \apprord x'$
        \end{itemize}
\end{itemize}

\begin{center}
    \includegraphics[width=8cm]{assets/order1}
\end{center}

\subsubsection{Liste e funzioni}
Per le liste si ha:
\begin{itemize}
    \item $\bot \apprord xs$
    \item $[ \ ] \apprord xs \iff xs = [ \ ] \land (x:xs) \not\apprord [ \ ]$
    \item $x:xs \apprord y:ys \iff x \apprord y \land xs \apprord ys$
\end{itemize}

Per esempio, $[1, \bot, 3]$ e $1:2:\bot$ sono approssimazioni di $[1,2,3]$.

\begin{center}
    \includegraphics[width=8cm]{assets/order2}
\end{center}

Invece, una funzione $f$ ne approssima un'altra $g$ se per ogni possibile input $x$, l'output di $f$ approssima quello di $g$:
$$\text{date } f,g: a \to b \text{ si ha } f \apprord_{a\to b} g \iff \forall x.fx \apprord_b gx$$

\subsubsection{Ordine e limiti}

\begin{defframe}{Complete Partial Order}{}
Se abbiamo una catena di approssimazioni $x_1 \apprord x_2 \apprord x_3 \dots$, questa tenderà ad un limite $\sqcup x = \lim_{n\to\infty} x_n$ che soddisfa le seguenti condizioni:
\begin{itemize}
    \item per ogni $n$, $x_n \apprord \sqcup x$ (il limite è un maggiorante)
    \item se per ogni $n$, $x_n \apprord y$, allora $\sqcup x \apprord y$ (il limite è il sup)
\end{itemize}
{\color{gray}(simile al lemma di Zorn ma più forte)}

Un ordinamento che soddisfa queste proprietà si definisce \textbf{Complete Partial Order}
\end{defframe}

I tipi per domini di funzioni computabili (e quindi eseguibili da una macchina) sono solitamente CPO.

\begin{defframe}{Funzioni computabili}{}
    Una funzione computabile gode di due proprietà:
    \vspace{-0.5em}
    \begin{enumerate}
        \item \textbf{monotonia}: $x \apprord y \implies f x \apprord f y$
        \item \textbf{continuità}: $f(\sqcup x) = \sqcup f x$
    \end{enumerate}
\end{defframe}

\begin{gframe}{Approfondimento: computabilità logica vs. semantica}
È interessante comparare la definizione di funzione computabile fornita dalla teoria dei domini a quella classica della logica matematica (funzioni $\mu$-ricorsive). 

Mentre la definizione logica è di tipo costruttivo - identifica la classe delle funzioni calcolabili come la minima classe ottenuta partendo da funzioni base (zero, successore, proiezioni) e applicando operatori di chiusura (composizione, minimalizzazione) -, la definizione basata sui CPO è di tipo semantico. Essa stabilisce i vincoli necessari affinché una funzione possa operare su informazioni parziali o infinite:

\begin{itemize}
    \item La \textit{monotonia} garantisce che la computazione sia un processo ad accumulo di informazione: se l'input diventa più definito ($x \apprord y$), l'output non può diminuire o cambiare l'informazione già prodotta ($f x \apprord f y$)
    \item La \textit{continuità} formalizza la natura finitaria della calcolabilità: il valore di una funzione su un input infinito (limite) deve essere determinato esclusivamente da ciò che la funzione produce sulle approssimazioni finite di quell'input ($f(\sqcup x) = \sqcup f x$)
\end{itemize}

In questo contesto, la parzialità delle funzioni logiche viene "internalizzata" nel dominio tramite l'elemento $\perp$ (bottom), che rappresenta il livello minimo di informazione o una computazione non terminata. Questo permette di trattare la non-terminazione non come un'assenza di valore, ma come un valore matematico ben definito su cui è possibile applicare i teoremi di punto fisso per dare un senso alla ricorsione.
\end{gframe}

Alcuni esempi di funzioni non computabili sono:
\begin{itemize}
    \item $
f(x) = \begin{cases}
    0 & x = \bot \\
    1 & x \neq \bot
\end{cases}
$
\begin{itemize}
    \item non è computabile perché non è monotona: $\bot \apprord x$, ma $0 \not\apprord 1$
\end{itemize}

\item $
f(xs) = \begin{cases}
    \bot & xs \text{ è finita o parziale}\\
    1 & xs \text{ è infinita}
\end{cases}
$
\begin{itemize}
    \item non è computabile perché non è continua: $\sqcup g \ x_n = \bot \neq 1 = g(\sqcup \ x_n)$
\end{itemize}
\end{itemize}
\section{Teoremi di punto fisso}

Nella semantica denotazionale, l'esecuzione di un programma ricorsivo è vista come la ricerca di un "punto fisso" di una funzione. Per formalizzare questo concetto, dobbiamo estendere la nozione di CPO a strutture algebriche più ricche.

\begin{defframe}{Reticolo Completo}{}
    Un \textbf{reticolo completo} (Complete Lattice) $(L, \apprord)$ è un insieme parzialmente ordinato in cui \textit{ogni} sottoinsieme $A \subseteq L$ possiede sia un estremo inferiore che un estremo superiore:
    \begin{itemize}
        \item Estremo superiore (Supremum): $\sqcup A = \min \{ x \mid \forall a \in A.\ x \ge a \}$
        \item Estremo inferiore (Infimum): $\sqcap A = \max \{ x \mid \forall a \in A.\ x \le a \}$
    \end{itemize}
\end{defframe}

Avere un reticolo completo ci permette di enunciare due teoremi fondamentali che garantiscono l'esistenza (e la calcolabilità) dei punti fissi per le funzioni su questi domini (un valore $x$ è un \textit{punto fisso} per una funzione $f$ se $f(x) = x$).

\begin{defframe}{Teorema di Knaster-Tarski}{}
    In un reticolo completo $(L, \apprord)$, \textbf{ogni funzione monotona} ammette un massimo e un minimo punto fisso.
\end{defframe}

Il teorema di Knaster-Tarski assicura l'esistenza del punto fisso, ma non dice come trovarlo. Il teorema di Kleene ci fornisce invece un metodo costruttivo.

\begin{defframe}{Teorema di Kleene (del minimo punto fisso)}{}
    Se una funzione $f$ è \textbf{continua}, il suo minimo punto fisso (Least Fixed Point, LFP) può essere costruito iterativamente partendo dal livello minimo di informazione ($\bot$).
    
    Esso è definito come il limite delle applicazioni ripetute di $f$:
    \[ \text{LFP}(f) = f^\omega(\bot) = \lim_{n\to\infty} f^n(\bot) = \sqcup \{ \bot, f(\bot), f(f(\bot)), \dots \} \]
    \begin{proofframe}
    Sappiamo che la sequenza $\bot \apprord f(\bot) \apprord f(f(\bot)) \dots$ è una catena monotona crescente e che quindi (poiché ci troviamo in un CPO) ha un estremo superiore.

    Possiamo dimostrare che questo estremo superiore $f^\omega$ è un punto fisso sfruttando la continuità:
    \begin{align*}
        f^\omega &= \sup \{ f^i(\bot) \mid i < \omega \} \quad &\text{(per definizione di } f^\omega \text{)} \\
        &= \sup \{ f(f^i(\bot)) \mid i < \omega \} \quad &\text{(spostando l'indice di 1 passo)} \\
        &= f(\sup \{ f^i(\bot) \mid i < \omega \}) \quad &\text{(\textbf{per la continuità di } } f\text{)} \\
        &= f(f^\omega) \quad &\text{(sostituendo la definizione di } f^\omega\text{)}
    \end{align*}
    Poiché $f(f^\omega) = f^\omega$, $f^\omega$ è per definizione un punto fisso di $f$.
    \end{proofframe}
\end{defframe}

\begin{gframe}{Significato computazionale dei teoremi di punto fisso}
Quando scriviamo una definizione ricorsiva come \code{ones = 1 : ones}, stiamo implicitamente definendo un'equazione: $x = f(x)$, dove la funzione è $f(x) = 1 : x$ - stiamo quindi cercando il punto fisso di $f$.

Il Teorema di Kleene ci spiega esattamente cosa fa la valutazione lazy: il computer costruisce la lista infinita partendo dal non-calcolato ($\bot$) e applicando la funzione in modo continuo:
\begin{enumerate}
    \item $f^0(\bot) = \bot$ \hfill \textit{(nessuna informazione)}
    \item $f^1(\bot) = 1 : \bot$ \hfill \textit{(primo elemento noto)}
    \item $f^2(\bot) = 1 : 1 : \bot$ \hfill \textit{(secondo elemento noto)}
\end{enumerate}
Il limite di questo processo all'infinito ($f^\omega$) è proprio la nostra lista infinita. La continuità garantisce che per estrarre una quantità di informazione finita dall'output (es. \code{take 2 ones}), il computer debba compiere solo un numero finito di passi, senza bloccarsi cercando di valutare l'intera struttura infinita.
\end{gframe}

\section{Induzione e co-induzione}

Nella teoria dei tipi e nella programmazione funzionale, le strutture dati possono essere definite in due modi duali: tramite induzione (costruendo valori) o tramite co-induzione (osservando valori). Queste due visioni corrispondono alla ricerca di minimi e massimi punti fissi.

\begin{defframe}{Insiemi (Algebre) Induttivi}{}
    Un insieme \textbf{induttivo} $A$ viene definito "partendo dal basso" (tipicamente dall'insieme vuoto), definendo le regole (costruttori) che permettono di creare valori di tipo $A$.
    
    \begin{itemize}
        \item $A$ è il \textbf{minimo} insieme chiuso rispetto alle operazioni di costruzione (rappresentato da un assioma di minimalità, es. "nient'altro è un numero naturale");
        \item L'insieme $A$ rappresenta il \textbf{minimo punto fisso} di un'operazione di costruzione;
        \item Gli elementi di $A$ sono il risultato di un'applicazione \textbf{finita} dei costruttori, a partire da costruttori costanti.
            \item Da un punto di vista logico, ad $A$ è associato un \textbf{principio di induzione} per provare proposizioni su $A$.
            \item Da un punto di vista informatico (minima algebra / oggetto iniziale in Teoria delle Categorie), ad $A$ è associato un principio di \textbf{ricorsione e/o iterazione} (come \code{fold}) per definire funzioni che "consumano" il tipo: $f \colon A \to B$.
    \end{itemize}
\end{defframe}

\begin{defframe}{Insiemi (Co-algebre) Co-induttivi}{}
    Un insieme \textbf{co-induttivo} $A$ viene definito "partendo dall'alto" (ad esempio da un universo $U$), definendo le regole (distruttori) che permettono di riconoscere i valori che \textit{non} appartengono ad $A$.
    
    \begin{itemize}
        \item $A$ è il \textit{massimo} insieme chiuso rispetto alle operazioni di eliminazione. 
        \item L'insieme $A$ rappresenta il \textbf{massimo punto fisso} di un'operazione di eliminazione.
        \item Gli elementi scartati ($U \setminus A$) vengono riconosciuti dopo un numero finito di applicazioni dei distruttori, ma per gli elementi di $A$ questo processo può continuare all'infinito (es. liste infinite o stream).
            \item Da un punto di vista logico, ad $A$ è associato un \textbf{principio di co-induzione} (come la bisimulazione), usato solitamente per provare \textit{equivalenze} su $A$.
            \item Da un punto di vista informatico (massima algebra / oggetto finale in Teoria delle Categorie), ad $A$ è associato un principio di \textbf{co-ricorsione e/o co-iterazione} (come \code{unfold} o \code{iterate}) per definire funzioni che "generano" o costruiscono oggetti di tipo $A$: $f \colon B \to A$.
    \end{itemize}
\end{defframe}

\subsection{Insiemi induttivi}
\begin{gframe}{Approfondimento: lemma di Lambek}
        Se consideriamo l'algebra induttiva come un \textit{oggetto iniziale} nella Teoria delle Categorie, il principio di ricorsione in programmazione è l'incarnazione computazionale dell'\textbf{omomorfismo unico} definito dal Lemma di Lambek (vedi \href{https://raw.githubusercontent.com/AglaiaNorza/notes-ig/main/terzo%20anno/linguaggi%20di%20programmazione.pdf}{\underline{appunti del corso ``Linguaggi di programmazione''} })
    
    Definire una funzione ricorsiva $f \colon A \to B$ su un tipo induttivo $A$ significa costruire l'unico omomorfismo unico che mappa i costruttori di $A$ nelle operazioni di un'altra algebra $B$.
\end{gframe}

    Ogni tipo induttivo ha il proprio operatore fondamentale (un costrutto di ordine superiore) che si occupa di "costruire" questo omomorfismo:
    \begin{itemize}
    \item \textbf{Tipi ``non ricorsivi'' (Booleani, Either):} anche se raramente definiti come tipi ricorsivi, possiedono un principio di ricorsione. Per i booleani è semplicemente l'\code{if-then-else}, mentre per il tipo \code{Either} è l'applicazione di due funzioni $g$ e $h$ a seconda del costruttore (\code{Left} o \code{Right}). Il principio di induzione associato è l'analisi per casi.
        
        \item \textbf{Numeri naturali ($\mathbb{N}$):} l'operatore matematico che costruisce l'omomorfismo è \textbf{$\rho$ (iteratore)}. In Haskell, esso prende la forma della funzione \code{iter}:
        \begin{haskellbox}{}
iter f b 0 = b
iter f b n = f (iter f b (n-1))
\end{haskellbox}

        Questo comportamento ricalca esattamente la semantica dei \textbf{Numerali di Church}. 

Per ottenere una potenza espressiva maggiore, si passa allo schema della \textit{ricorsione primitiva} (l'operatore \code{rec}), in cui la funzione $f$ riceve in input anche il contatore $(n-1)$.
        
        \item \textbf{Liste finite:} sulle liste l'operazione che costruisce l'omomorfismo "sostituendo" i costruttori della lista con operazioni su un altro tipo è la \code{foldr}. Rappresenta il principio di iterazione/ricorsione generalizzato per le liste:
    \begin{haskellbox}{}
iter f b [] = b
iter f b (x:xs) = f x (iter f b xs)  -- corrisponde a foldr f b xs
    \end{haskellbox}
    \end{itemize}

\subsection{Insiemi coinduttivi}
L'insieme coinduttivo che origina dai costruttori dei numeri naturali è formato dai naturali, più un \textbf{naturale infinito non-standard} $\omega = S^{\infty}$, ottenuto da infinite applicazioni del costruttore $S$.

L'insieme coinduttivo più interessante è quello delle liste infinite (o \textbf{Stream}). Non ha senso definire funzioni per ricorsione su di esse, ma si possono sempre definire funzioni che creano Stream di cui, in tempo finito, si può accedere ad una porzione finita.

Il principio di ricorsione sulle Stream si può definire tramite due funzionali:
\begin{itemize}
    \item \code{iterate}, che permette di definire una funzione $f: A \to Stream \ B$
        \subitem data una funzione \code{f :: a -> b } ed un valore \code{x :: b}, permette di costruire uno stream in questo modo:
        \begin{haskellbox}{}
iterate f x = x : iterate f (f x)
        \end{haskellbox}

        (crea quindi lo stream \code{x : f x : f (f x) : ... })
\end{itemize}

Conviene però avere un principio di ricorsione più generale: la cosiddetta \textbf{co-ricorsione}, ottenibile con il funzionale \code{unfold}.

\begin{haskellbox}{}
unfold :: (b -> (a, b)) -> b -> Stream a
unfold f y = x : unfold f y' where (x, y') = f y
\end{haskellbox}

Se invece vogliamo considerare l'algebra che contiene la lista vuota, allora la massima co-algebra contiene anche le liste finite. In tal caso, possiamo definire il funzionale \code{unfoldLS}, che costruisce anche liste finite:
\begin{haskellbox}{}
unfold :: (b -> Bool) -> (b -> (a, b)) -> b -> [a]
unfold p f y = if p y then []
            else x : unfold p f y' where (x, y') = f y
\end{haskellbox}

\subsection{Bisimulazione}
Come per gli insiemi induttivi, abbiamo bisogno di un modo per derivare equivalenze su insiemi co-iduttivi. La controparte dell'induzione che ci permette di farlo è la \textbf{bisimulazione}.

L'idea alla base della bisimulazione è osservare se due processi (o liste) ``si comportano'' allo stesso modo per ogni osservazione finita che possiamo farne.

Si procede in questo modo:
\begin{enumerate}
    \item si definisce una relazione $R$ tra due liste
    \item se si dimostra che per ogni coppia $(xs, ys) \in R$, le loro teste sono uguali e le loro code sono ancora in $R$, allora le liste si dicono \textbf{bisimili} ($\approx$)
\end{enumerate}

Formalmente, si definisce una bisimulazione $R$ tra coppie di liste come segue:
\begin{align*}
    R \ xs \ ys \implies & xs = ys = \bot \text{ oppure } xs = ys = [ \ ] \\
                         & \text{oppure } \exists v, vs, ws \text{ t.c. } xs = v:vs \land ys = v:ws \land vs \ R \ ws
\end{align*}

Definiamo quindi $xs \approx ys$ se esiste una bisimulazione $R$ tale che $xs \ R \ ys$, il che equivale a scegliere la \textbf{massima equivalenza}, definita come massimo punto fisso.

\begin{gframe}{Prova per bisimulazione}
    Possiamo dimostrare per bisimulazione che:
    \begin{haskellbox}{}
map f (iterate f x) = iterate f (f x)
    \end{haskellbox}
    Ci basta dimostrare che la relazione $R$ che contiene:
    $$\{ \code{map f (iterate f x), iterate f (f x)} \mid \code{f,x} \text{ opportuni} \}$$
    è una bisimulazione.

    Facendo ricorso all'algebra dei programmi, possiamo ``svolgere'' le due espressioni:
    \begin{haskellbox}{}
map f (iterate f x) = map f (x : iterate f (f x))
                    = f x : map f (iterate f (f x))

    \end{haskellbox}
e anche
    \begin{haskellbox}{}
iterate f (f x) = f x : iterate f (f (f x))
    \end{haskellbox}

    Entrambe producono la stessa testa (\code{f x}), e le code sono in $R$ (in cui la nuova funzione applicata sarà \code{f.f} invece di \code{f}).
\end{gframe}

\subsection{Dimostrazione per approssimazione}
L'alternativa alla co-induzione consiste nell'effettuare una dimostrazione per induzione sulle approssimazioni finite di una lista. Per farlo, si internalizza in Haskell la funzione \code{approx}, strettamente imparentata con \code{take}:

\begin{haskellbox}{}
approx 0 xs = undefined
approx n [] = []
approx n (x:xs) = x : approx (n-1) xs
\end{haskellbox}

La proprietà fondamentale di \code{approx} è che il limite delle approssimazioni di una lista è la lista stessa:
$$\sqcup_{n} \text{approx } n \ xs = xs$$

Di conseguenza, per dimostrare l'uguaglianza tra due liste $xs$ e $ys$, è sufficiente mostrare che tutte le loro approssimazioni finite siano uguali per ogni $n$:
$$\forall n, \text{approx } n \ xs = \text{approx } n \ ys \implies xs = ys$$

\begin{gframe}{Esempio}
Proviamo a dimostrare nuovamente la proprietà:
$$\text{approx } n \ (\text{map } f \ (\text{iterate } f \ x)) = \text{approx } n \ (\text{iterate } f \ (f \ x))$$

\begin{itemize}
    \item \textbf{caso base:} per $n = 0$ la proprietà è banalmente vera perché \code{approx 0 xs} restituisce sempre $\perp$ (undefined).

    \item \textbf{caso induttivo:} si procede riducendo le due espressioni separatamente tramite program calculation e applicando l'ipotesi induttiva:

\begin{align*}
    & \text{approx } (n+1) \ (\text{iterate } f \ (f \ x)) \\
    = & \{ \text{def. di iterate} \} \\
    & \text{approx } (n+1) \ (f \ x : \text{iterate } f \ (f \ (f \ x))) \\
    = & \{ \text{def. di approx} \} \\
    & f \ x : \text{approx } n \ (\text{iterate } f \ (f \ (f \ x))) \\
    = & \{ \text{ipotesi induttiva} \} \\
    & f \ x : \text{approx } n \ (\text{map } f \ (\text{iterate } f \ (f \ x))) \\
    = & \{ \text{def. di approx al contrario} \} \\
    & \text{approx } (n+1) \ (f \ x : \text{map } f \ (\text{iterate } f \ (f \ x))) \\
    = & \{ \text{def. di map al contrario} \} \\
    & \text{approx } (n+1) \ (\text{map } f \ (x : \text{iterate } f \ (f \ x))) \\
    = & \{ \text{def. di iterate al contrario} \} \\
    & \text{approx } (n+1) \ (\text{map } f \ (\text{iterate } f \ x))
\end{align*}
\end{itemize}
\end{gframe}

\section{Funtori e Applicativi}

\begin{gframe}{Concetti principali}
    Sia i Funtori che gli Applicativi servono a gestire computazioni all'interno di un \textbf{contesto}, ma con ``poteri'' diversi:
    \begin{itemize}
        \item \textbf{Functors}: permettono di applicare una funzione "pura" a un valore nel contesto - modificano il contenuto ma non la struttura del contesto;
        \item \textbf{Applicatives}: permettono di applicare funzioni che sono esse stesse dentro un contesto a valori nel contesto - sono necessari per gestire funzioni con più argomenti;
    \end{itemize}
\end{gframe}

\begin{center}
\begin{tabular}{|l|l|l|}
\hline
\textbf{classe} & \textbf{operatore} & \textbf{scopo} \\ \hline
\code{Functor} & \code{fmap :: (a -> b) -> f a -> f b} & trasformazione (1 argomento) \\ \hline
\code{Applicative} & \code{<*> :: f (a -> b) -> f a -> f b} & combinazione ($n$ argomenti) \\ \hline
\end{tabular}
\end{center}

\subsection{Funtori}

L'idea alla base di un \textbf{funtore} è la generalizzazione del funzionale \code{map}, tipicamente utilizzato per le liste, a tutti i possibili costruttori di tipo.

Se si ha:
\begin{itemize}
    \item una funzione $f:: a\to b$
    \item un costruttore di tipo $T$ che trasforma un tipo $a$ in un tipo $T a$ (come per esempio \code{[a]} oppure \code{Maybe a})
\end{itemize}

Si può ottenere in modo canonico una funzione
$$T f :: T a \to T b$$

\begin{center}
\begin{tikzpicture}[x=5cm, y=2.5cm, >=latex, very thick, font=\Large]
    \node (a)  at (0,0) {$a$};
    \node (b)  at (1,0) {$b$};
    \node (Ta) at (0,1) {$Ta$};
    \node (Tb) at (1,1) {$Tb$};
    
    \node[text=teal] at (0.5, 0.5) {\huge $\mathbf{=}$};
    
    \draw[->] (a) -- node[below=4pt] {$f$} (b);
    \draw[->] (a) -- (Ta);
    \draw[->] (b) -- (Tb);
    
    \draw[->, teal] (Ta) -- node[above=4pt, text=green!60!black] {$Tf$} (Tb);
\end{tikzpicture}
\end{center}

\begin{gframe}{}
Nel mondo categoriale, il concetto di ``funtore'' descrive un passaggio tra due categorie: dati due oggetti (tipi come $a$ o $b$), un funtore agisce come un costruttore di tipo ($T$) - prende un oggetto $a$ nella categoria $C$ e lo trasforma in un oggetto $T a$ nella categoria $D$.
\begin{itemize}
    \item il punto importante è che un funtore non mappa solo i tipi, ma anche le relazioni tra essi: se esiste una funzione $f : a \to b$, il funtore deve fornire un modo per ottenere una funzione corrispondente $T f$ che operi sui tipi $T a \to T b$
    \item un funtore è quindi una \textbf{mappa che preserva la struttura}
\end{itemize}
\end{gframe}

La classe \code{Functor} richiede che sia definita una funzione \code{fmap}, che rende commutativo il passaggio dai tipi base ai tipi ``inscatolati''

\begin{haskellbox}{}
class Functor t where
    fmap :: (a -> b) -> t a -> t b
\end{haskellbox}
\begin{itemize}
    \item il tipo di \code{fmap} è parametrico, oltre che nei tipi \code{a} e \code{b}, anche rispetto a un \textbf{costruttore di tipo} \code{t} (si vede che \code{t} deve essere un costruttore di tipo perché si applica ad altri tipi)
\end{itemize}

\begin{gframe}{Esempi}
    Alcuni esempi noti di istanze di funtori sono:
    \begin{itemize}
        \item le liste
\begin{haskellbox}{}
class Functor [] where
    fmap = map
\end{haskellbox}

\item il tipo \code{Maybe}
\begin{haskellbox}{}
class Functor Maybe where
    fmap f Nothing = Nothing
    fmap f (Just x) = Just (f x)
\end{haskellbox}

\item gli alberi binari
\begin{haskellbox}{}
class Functor BinTree where
    fmap f (F x) = F (f x)
    fmap f (R r lft rgt) = 
        R (f r) (fmap f lft) (fmap f rgt)
\end{haskellbox}
    \end{itemize}
\end{gframe}

Definire un tipo come istanza di \code{Functor} permette di scrivere codice generico che funziona su diverse strutture dato. 

Per esempio, si può definire una funzione \code{inc} che incrementa i valori all'interno di qualsiasi funtore:
\begin{haskellbox}{}
inc :: Functor t => t Int -> t Int
inc = fmap (+1)
\end{haskellbox}

\subsubsection{Functor laws}
Affinché un funtore sia definito correttamente, è necessario che rispetti le due functor laws:
\begin{enumerate}
    \item \code{fmap id = id}
    \item \code{fmap (f . g) = fmap f . fmap g}
\end{enumerate}

Nota bene: queste leggi non possono essere verificate dal type-checker, ma sono responsabilità del programmatore Haskell !

\begin{gframe}{Funtori come omomorfismi}
    Possiamo notare, osservando le \textit{functor laws}, che un funtore è effettivamente un \textbf{omomorfismo tra categorie} (un omomorfismo ``classico'' preserva l'operazione di gruppo, mentre un funtore preserva la ``struttura'' di una categoria, ovvero gli oggetti - i tipi -, e i morfismi - le funzioni.)
\end{gframe}

\subsection{Applicativi}
Un normale funtore permette di usare \code{fmap} con un solo argomento, ma è naturale voler generalizzare: si vuole considerare una forma astratta di applicazione, che generalizzi l'usuale applicazione di funzioni. Per esempio, si vuole poter propagare un fallimento nel caso di \code{Maybe}, o applicare una funzione $n$-aria a valori ``inscatolati'' senza dover istanziare $n$ funtori.

Haskell definisce quindi la classe \code{Applicative}.

\begin{haskellbox}{}
class Functor t => Applicative t where
    pure :: a -> t a
    (<*>) :: t (a -> b) -> t a -> t b
\end{haskellbox}

Un applicativo estende un funtore, e richiede due operatori fondamentali:
\begin{itemize}
    \item \code{pure}: prende un valore ``normale'' e lo immerge nel contesto (lo ``impacchetta'')
        \subitem per esempio, per \code{Maybe}, \code{pure x} è \code{Just x}
    \item \code{<*>} (\textit{splat} o \textit{apply}): prende una funzione ``impacchettata'' in un contesto (\code{t (a -> b)}) e un valore anch'esso in un contesto (\code{t a}), ``scarta'' concettualmente entrambi i contesti, applica la funzione al valore (gestendo gli effetti del contesto), produce il risultato e lo ``incarta'' nel contesto (\code{t b})
\end{itemize}

\begin{gframe}{Esempio: Maybe}
    Possiamo istanziare \code{Maybe} anche come applicativo:
    \begin{haskellbox}{}
instance Applicative Maybe where
    pure = Just
    Nothing <*> _ = Nothing -- se uno dei due è Nothing, entrambi sono Nothing
    (Just g) <*> mx = fmap g mx
    -- <*> prende una funzione e un valore in un contesto, quindi
    -- per es. Just (+3) <*> Just 2 = Just 5
    \end{haskellbox}
\end{gframe}

\subsubsection{Liste come applicativi}
L'istanza del Prelude delle liste non si comporta come una \code{zipWith} nell'applicare le funzioni presenti in una lista, ma modella un non-determinismo. Una lista non viene vista come una sequenza, ma come un valore che può assumere più risultati posssibili contemporaneamente. Quindi, quando si applica una lista di funzioni ad una lista di valori, l'applicativo contiene tutte le combinazioni possibili (prodotto cartesiano).

\begin{haskellbox}{}
instance Applicative [] where
    -- pure : a -> [a]
    pure x = [x]

    -- <*> : [a -> b] -> [a] -> [b]
    gs <*> xs = [g x | g <- gs, x <- xs]
    -- ^ prodotto cartesiano 
\end{haskellbox}

Se eseguiamo quindi \code{pure (*) <*> [1,2] <*> [3,4]}:
\begin{itemize}
    \item \code{pure (*)} crea \code{[(*)]}
    \item \code{[(*)] <*> [1,2]} produce una lista di funzioni \code{[(1*), (2*)]}
    \item \code{[(1*), (2*)] <*> [3,4]} applica entrambe le funzioni ad entrambi i numeri
        \subitem si ottiene quindi \code{1*3, 1*4, 2*3, 2*4} $\to$ \code{[3, 4, 5, 6]}
\end{itemize}

A volte si vuole però che le liste si comportino come \code{zipWith} - per questo, Haskell introduce \code{ZipList}

\begin{haskellbox}{}
 -- per poter definire un'altra istanza di Applicative senza conflitti
newtype ZipList a = Z [a]

instance Applicative ZipList where
    -- pure : a -> ZipList a
    pure x = Z (repeat x) -- crea una lista infinita di x

    -- <*> : [a -> b] -> [a] -> [b]
    Z fs <*> Z xs = Z (zipWith (\f x -> f x) fs xs)
\end{haskellbox}

\begin{itemize}
    \item \code{pure} è \code{repeat} così che si adatti a qualsiasi lunghezza della lista di valori con cui verrà  accoppiata (se restituisse un solo elemento, \code{zipWith} si fermerebbe subito)
    \item \code{<*>} è \code{zipWith} - applica la funzione in posizione 1 al valore in posizione 1 ecc
\end{itemize}

\subsubsection{Applicative laws}
Come i funtori, anche gli applicatori hanno una serie di regole da rispettare:
\begin{enumerate}
    \item \code{<*>} \textbf{preserva l'identità} 
        \subitem \code{pure id <*> x = x}
        \subitem se prendiamo la funzione identità, la inseriamo in un contesto e la applichiamo ad un valore \code{x} già contestualizzato, il risultato deve rimanere \code{x}
    \item \code{<*>} \textbf{preserva l'applicazione di funzioni}
        \subitem \code{pure (g x) = pure g <*> pure x}
        \subitem applicare \code{g} ad un valore \code{x} e poi ``impacchettare'' il risultato dà lo stesso esito che applicare la versione ``impacchettata'' di \code{g} alla versione ``impacchettata'' di \code{x}
    \item \textbf{non conta l'ordine} con cui si fa embedding nel contesto
        \subitem \code{x <*> pure y = pure (\textbackslash g -> g y) <*> x}
    \item \textbf{composizione} di \code{<*>}
        \subitem \code{x <*> (y <*> z) = (pure (.) <*> x <*> y) <*> z}
        \subitem (forma di associatività per \code{<*>})
        \subitem la composizione di funzioni \code{(.)} può essere ``lifted'' a livello applicativo per garantire che la sequenza delle applicazioni sia coerente
\begin{itemize}
    \item ovvero, applicare due funzioni una dopo l'altra è identico a comporle prima in un'unica funzione e poi applicarla:
        \subitem \code{x <*> (y <*> z)}
        \subitem prima si applica la funzione \code{y} al valore \code{z} (ottenendo un risultato nel contesto), e poi si applica \code{x} a tale risultato - è quindi \code{x(y(z))}
    \subitem \code{(pure (.) <*> x <*> y) <*> z}
\subitem qui l'operatore di composizione standard \code{(.)} viene ``sollevato'' (lifted) nel contesto, permettendo di fondere le due funzioni \code{x} e \code{y} dentro il contesto - si crea quindi una funzione composta che sarà applicata a \code{z}
\subitem quindi, si applica \code{x . y} a \code{z}, ottenendo anche qui \code{x(y(z))}
\end{itemize}
\end{enumerate}

\section{Monadi Base}
Le monadi rappresentano il passo successivo rispetto agli applicativi per la gestione controllata di \textbf{side-effects}. Se un applicativo permette di combinare calcoli indipendenti, la monade permette di gestire calcoli in cui il passo successivo \textbf{dipende} dal risultato di quello precedente.

\subsection{Introduzione}
Per comprendere perché le monadi siano utili, consideriamo un valutatore di espressioni aritmetiche che supporta costanti e divisioni:

\begin{haskellbox}{}
data Term = Const Int | Div Term Term

eval :: Term -> Int
eval (Const a) = a
eval (Div t u) = eval t `div` eval u
\end{haskellbox}

Questa implementazione è elegante ma fragile: un'espressione come \code{Div (Const 1) (Const 0)} farà crashare il programma.

Si può usare il tipo \code{Maybe Int} per modellare il possibile fallimento. Tuttavia, il codice diventa immediatamente prolisso a causa della propagazione manuale dell'errore:

\begin{haskellbox}{}
eval :: Term -> Maybe Int
eval (Const a) = Just a
eval (Div t u) = case eval t of
                   Nothing -> Nothing
                   Just n  -> case eval u of
                                Nothing -> Nothing
                                Just m  -> if m == 0 then Nothing 
                                           else Just (n `div` m)
\end{haskellbox}

La divisione (il vero senso della funzione) è offuscata dalla gestione dei \code{Nothing}.

Si potrebbe pensare di usare l'operatore \code{<*>} degli applicativi per pulire il codice. Tuttavia, gli applicativi sono insufficienti quando la struttura del calcolo dipende dal valore dei risultati intermedi:

\begin{itemize}
    \item L'operatore \code{<*>} (\code{f (a -> b) -> f a -> f b}) richiede che la funzione da applicare sia "pura" e che i contesti (\code{f}) siano già definiti. In un calcolo applicativo, l'intera sequenza di effetti è determinata a priori: non si può usare il valore estratto da un \code{Maybe} per decidere quale sarà il prossimo "effetto".
    
    \item Se provassimo a usare una funzione che genera essa stessa un contesto, come 

        \code{safeDiv :: Int -> Int -> Maybe Int}, l'applicazione risulterebbe in un tipo

        \code{Maybe (Maybe Int)}. 
    \begin{itemize}
        \item e.g. \code{pure safeDiv <*> Just 10 <*> Just 0} produce \code{Just (Nothing)}
    \end{itemize}
    L'applicativo non possiede un meccanismo per ``appiattire'' questi due livelli di contesto in uno solo.
    
\end{itemize}
 Nella divisione, il fatto che il secondo argomento sia zero determina se il calcolo successivo deve fallire o meno. In Haskell, questa capacità di far sì che il ``guscio'' esterno (il contesto) venga generato dal valore contenuto nel passo precedente è l'essenza delle Monadi.

    In sintesi: gli Applicativi modellano calcoli indipendenti (i contesti sono paralleli), mentre le Monadi (come ora vedremo) modellano calcoli sequenziali, dove il risultato di un'operazione influenza la creazione della successiva.

\begin{gframe}{}
In termini di Teoria delle Categorie, l'applicativo gestisce contesti indipendenti (funtori monoidali), mentre per concatenare funzioni del tipo $A \to M(B)$ (chiamate frecce di Kleisli) abbiamo bisogno di una Monade.
\end{gframe}

\subsection{Monadi (shall I be pure or impure?)}
Molti aspetti della programmazione (come strutture dati mutabili e input/output) non hanno una rappresentazione efficace nel mondo funzionale, mentre altri potrebbero beneficiare da una programmazione con side-effects (come eccezioni e tracciamento delle computazioni).

Alcuni linguaggi funzionali, come ML, introducono costrutti non funzionali come puntatori, eccezioni, ecc.; Haskell cerca un'integrazione basata su concetti matematici.

La classe \code{Monad} si può definire tramite gli applicativi in questo modo:
\begin{haskellbox}{}
class Applicative m => Monad m where
    (>>=)  :: m a -> (a -> m b) -> m b  -- operatore "bind"
    (>>)   :: m a -> m b -> m b -- "sequence"/"then"
    return :: a -> m a  -- identico a "pure"
\end{haskellbox}
\begin{itemize}
    \item \code{return} (o \code{pure}) prende un valore e lo inserisce nel contesto minimo necessario
    \item \code{>>=} (\textbf{bind}) prende un valore ``impacchettato'' (\code{m a}), lo ``spacchetta'', e lo passa ad una funzione che accetta un valore puro  e restituisce un nuovo valore impacchettato (\code{a -> mb})
        \begin{itemize}
            \item in questo modo, la computazione viene sequenzializzata: prima viene valutato \code{m : M a} e poi il risultato viene passato a \code{f}
        \end{itemize}
    \item \code{>>} è una versione semplificata del bind: mentre \code{>>=} prende il valore della prima computazione e lo passa ad una funzione, \code{>>} esegue la computazione, ne scarta il valore prodotto e passa direttamente alla seconda computazione (il suo tipo è infatti quello della seconda computazione - in una computazione sequenziale, l'ultima istruzione è quella che determina il valore finale)
        \begin{itemize}
            \item \code{m >> f} è equivalente a \code{m >>= \_ -> f}
            \item serve quindi quando si vogliono solo i side-effects ma non quello che viene ritornato (e.g. operazioni di scrittura)
            \item in ogni caso, anche se il valore di tipo \code{a} viene scartato, il contesto (effetto computazionale) di \code{m a} viene preservato e combinato con quello di \code{b}
        \end{itemize}
\end{itemize}

In questo modo, una funzione \code{f : a -> b} viene rimpiazzata da \code{f' : a -> M b}, in cui \code{M} \textbf{cattura gli effetti computazionali} di \code{f}.

\begin{gframe}{Valutatore generalizzato}
    Possiamo quindi scrivere una versione generale del nostro valutatore che accetti qualsiasi monade:
\begin{haskellbox}{}
eval :: Monad m => Term -> m Int
eval (Const a) = pure a
eval (Div t u) = eval t >>= \a ->
                 eval u >>= \b ->
                 pure (a `div` b)
\end{haskellbox}
\end{gframe}

Segue un link utile per comprendere al meglio questi concetti: \href{https://www.adit.io/posts/2013-04-17-functors,_applicatives,_and_monads_in_pictures.html}{[qui]}

\subsection{Alcune monadi note}

\subsubsection{Monade identità}
La \textbf{monade identità} è simile all'elemento neutro di un'algebra. Serve a rappresentare una computazione senza un effetto collaterale, e permette di definire una funzione generale ed usarla senza aggiungere nessun comportamento ``speciale''.

\begin{haskellbox}{}
instance Identity Monad where
    return a = a
  -- (>>=)  :: Identity a -> (a -> Identity b) -> Identity b
    a >>= f = f a
\end{haskellbox}

\subsubsection{Liste}
Anche le \textbf{liste} si possono definire come monadi.
\begin{haskellbox}{}
instance [] Monad where
    return x = [x]
    xs >>= f = concat (map f xs)
    -- xs >>= f = [y | x <- xs, y <- f x]
\end{haskellbox}

La logica è quella di prendere una lista, applicare ad ogni elemento una funzione che restituisce a sua volta una lista, e poi ``appiattire'' tutto (\code{>>=} è analogo alle list comprehension).

\subsubsection{Monade eccezione}
Il tipo \code{Maybe} può essere trattato come una monade per gestire le eccezioni.

L'operatore \code {>>=} si comporta in questo modo:
\begin{itemize}
    \item se il risultato precedente è \code{Nothing}, questo sarà propagato
    \item se è \code{Just x}, \code{x} viene spacchettato e passato alla funzione successiva
\end{itemize}

\begin{haskellbox}{}
instance Maybe Monad where
    return x = Just x
    Nothing >>= _ = Nothing
    (Just x) >>= f = f x
\end{haskellbox}

\subsection{Monade State Transformer}
In un mondo imperativo, una variabile può cambiare valore nel tempo. In Haskell,  i dati sono immutabili. Per simulare uno ``stato'' che cambia, si usa una funzione che prende lo stato corrente e restituisce sia un risultato che un nuovo stato.

Per definire una State Transformer ``manualmente'', servono:
\begin{itemize}
    \item uno stato - rappresenta l'informazione che evolve nel tempo (es. un contatore) - per esempio un intero
        \subitem \code{type State = Int}.
    \item uno State Transformer (ST) - un wrapper che racchiude la logica del cambiamento di stato. Per definirlo manualmente, si deve usare un costruttore fittizio \code{S} perché in Haskell non è possibile definire un sinonimo di tipo (\code{type}) come istanza di una classe:
\begin{haskellbox}{}
newtype ST a = S (State -> (a, State))
\end{haskellbox}
\end{itemize}

A causa di questo costruttore, è utile definire un'operazione di \textbf{applicazione} (\code{app}) che rimuove l'etichetta \code{S} e applica la funzione interna allo stato fornito.

\begin{haskellbox}{}
app :: ST a -> State -> (a, State)
app (S st) x = st x
\end{haskellbox}

Possiamo simulare effetti collaterali (come incrementare un contatore) definendo funzioni specifiche come \code{tick}. Tuttavia, senza l'uso delle monadi, dovremmo gestire manualmente il ``threading'' dello stato (passare il nuovo stato \code{s'} alla computazione successiva), rendendo il codice ripetitivo.

\begin{haskellbox}{}
-- "side effect" manuale
tick :: ST Int -> ST Int
tick (S st) = S (\s -> 
    let (a, s') = st s -- esegue la computazione e ottiene il nuovo stato
    in (a, s' + 1))    -- incrementa manualmente lo stato
\end{haskellbox}

Definendo \code{ST} come \textbf{monade}, l'operatore \code{>>=} si occupa di ``nascondere'' il passaggio dello stato aggiornato tra le funzioni.

\begin{haskellbox}{}
instance Monad ST where
    -- return mette un valore in uno stato neutro
    return x = S (\s -> (x, s)) 

    -- (>>=) concatena le azioni e fa fluire lo stato s' automaticamente
    st >>= f = S (\s -> 
        let (x, s') = app st s 
        in app (f x) s') 
\end{haskellbox}

In Haskell, le monadi sono strutturate secondo una gerarchia di classi. Una monade è, per definizione, un applicativo "potenziato". Pertanto, per definire \code{ST} come monade, dobbiamo obbligatoriamente istanziarla come applicativo, e di conseguenza anche come funtore.

Infatti:
\begin{enumerate}
    \item Functor permette di usare \code{fmap} per trasformare il risultato di una computazione di stato;
    \item Applicative permette di gestire computazioni indipendenti in sequenza e fornisce \code{pure};
    \item Monad introduce l'operatore \code{>>=} (bind), permettendo a una computazione di dipendere dal risultato della precedente;
\end{enumerate}

\begin{haskellbox}{}
instance Functor ST where
    -- fmap :: (a -> b) -> ST a -> ST b
    fmap f (S st) = S (\s -> 
        let (x, s') = st s 
        in (f x, s'))
\end{haskellbox}

\begin{haskellbox}{}
instance Applicative ST where
    pure x = S (\s -> (x, s))

    -- (<*>) :: ST (a -> b) -> ST a -> ST b
    stf <*> stx = S (\s ->
        let (f, s')  = app stf s in
        let (x, s'') = app stx s' in
        (f x, s''))
\end{haskellbox}

\begin{gframe}{Flusso di stati e dati}
Per comprendere come variano le operazioni, possiamo visualizzare il trasformatore di stato come una scatola nera con un ingresso (lo stato iniziale $s$) e due uscite (il valore prodotto $x$ e lo stato aggiornato $s'$).

\begin{itemize}
    \item \textbf{Pure}: il valore $x$ viene semplicemente "iniettato" nel sistema. Lo stato $s$ entra ed esce esattamente uguale, senza essere toccato - non c'è alcuna trasformazione, solo un trasporto di informazioni.
    
    \item \textbf{Funtore (fmap)}: lo stato $s$ entra nella scatola $st$ e produce un valore intermedio $x$ e uno stato $s'$. La funzione $g$ viene applicata solo al valore $x$, ottenendo $g(x)$. Lo stato $s'$ non viene alterato dalla funzione.

    \item \textbf{Applicativo (\code{<*>})}: rappresenta il sequenziamento di due trasformazioni indipendenti;
    \begin{enumerate}
        \item lo stato iniziale $s$ entra nel primo blocco, che produce una funzione $f$ e uno stato intermedio $s'$
        \item questo stato $s'$ viene passato come input al secondo blocco, che produce un valore $x$ e lo stato finale $s''$
        \item solo alla fine la funzione $f$ e il valore $x$ vengono combinati
    \end{enumerate}
    Qui lo stato fluisce "in serie" tra le due scatole, ma la natura della seconda operazione non dipende dal risultato della prima.

\item \textbf{Monade (\code{>>=})}: come nell'applicativo, lo stato fluisce in serie ($s \to s'$), ma il valore $x$ prodotto dalla prima scatola viene usato per scegliere o costruire la scatola successiva.
\end{itemize}

\end{gframe}

\subsection{Do notation}
La \textbf{do notation} è una sintassi alternativa (zucchero sintattico) per scrivere computazioni monadiche in modo leggibile, simulando uno stile imperativo. Ogni blocco \code{do} viene tradotto dal compilatore in una serie di applicazioni degli operatori \code{>>=} e \code{>>}.

Sia $p$ un'azione (un'espressione di tipo \code{Monad m => m a}) e \code{stmts} una sequenza non vuota di comandi. Valgono le seguenti uguaglianze:
\begin{itemize}
    \item \textbf{azione singola:} un blocco \code{do} contenente una sola azione coincide con l'azione stessa.
    \begin{haskellbox}{}
do {p} = p
    \end{haskellbox}

    \item \textbf{sequenziamento:} se un'azione è seguita da altri comandi, viene usato l'operatore \code{>>}; il risultato della prima azione viene scartato.
    \begin{haskellbox}{}
do {p; stmts} = p >> do {stmts}
    \end{haskellbox}

    \item \textbf{binding:} il comando \code{x <- p} estrae il valore dal contesto monadico di $p$ e lo rende disponibile per i comandi successivi tramite una funzione lambda.
    \begin{haskellbox}{}
do {x <- p; stmts} = p >>= \x -> do {stmts}
    \end{haskellbox}
\end{itemize}

\textbf{Nota bene:} un blocco \code{do} non può terminare con un comando di estrazione (\code{x <- p}), ma deve obbligatoriamente terminare con un'azione pura o monadica. Ad esempio, \code{do {x <- getChar}} non è un'espressione valida.


\subsubsection{Esempio Pratico: Tree Relabeling}
Per comprendere meglio il funzionamento pratico della monade \code{ST}, analizziamo il problema del Relabeling di un albero binario. L'obiettivo è prendere un albero ed etichettare ogni sua foglia progressivamente con un numero intero.

Definiamo il tipo di dato per l'albero: i dati risiedono nelle foglie (\code{F}), mentre i nodi di ramificazione (\code{R}) servono solo a biforcare la struttura:
\begin{haskellbox}{}
data Tree a = F a | R (Tree a) (Tree a)
\end{haskellbox}

Per generare gli indici progressivi, definiamo un'azione nello stato chiamata \code{fresh}, il cui scopo è leggere il contatore corrente e incrementarlo:
\begin{haskellbox}{}
fresh :: ST Int
fresh = S (\n -> (n, n + 1))
\end{haskellbox}

\begin{gframe}{Come funziona ST?}
Ricordando che il tipo \code{ST a} racchiude una funzione di transizione \code{State -> (a, State)}, quando viene invocata l'azione \code{fresh} accade questo:
\begin{itemize}
    \item la funzione accetta lo stato corrente \code{n} (l'intero che fa da contatore);
    \item restituisce la coppia \code{(valore\_prodotto, nuovo\_stato)}
        \subitem (che nello specifico è \code{(n, n + 1)})
    \item il valore prodotto \code{n} rappresenta il contatore prima dell'incremento, che verrà estratto ed assegnato alla foglia corrente;
    \item il \code{nuovo\_stato} (\code{n + 1}) non viene memorizzato nella variabile estratta, ma viene preso da Haskell per essere passato implicitamente alla riga di codice successiva;
\end{itemize}
   
\end{gframe}

Sfruttando la \code{do} notation, l'algoritmo di rietichettatura si scrive in modo sequenziale, mimando la sintassi di un linguaggio imperativo:
\begin{haskellbox}{}
mlabel :: Tree a -> ST (Tree Int)
mlabel (F _) = do 
  n <- fresh
  return (F n)

mlabel (R l r) = do 
  l' <- mlabel l
  r' <- mlabel r
  return (R l' r')
\end{haskellbox}

\begin{gframe}{Do notation}
Il compilatore Haskell converte il blocco \code{do} rimuovendo lo zucchero sintattico e traducendolo nell'operatore \code{>>=}:

\begin{haskellbox}{}
mlabel (F _) = fresh >>= \n -> return (F n)
\end{haskellbox}

Se sostituiamo l'operatore \code{>>=} con la sua definizione formale implementata per la monade \code{ST}, l'espressione diventa:
\begin{haskellbox}{}
mlabel (F _) = S (\s -> 
    let (n, s') = app fresh s 
    in app (return (F n)) s')
\end{haskellbox}

Quindi la computazione procede così:
\begin{enumerate}
    \item la computazione \code{fresh} viene spacchettata tramite \code{app} e applicata allo stato iniziale \code{s}; produce il valore di contatore corrente \code{n} (pari a \code{s}) e lo stato incrementato \code{s'} (pari a \code{s + 1}).
    \item l'operatore \code{>>=} cattura questa tupla, estrae il valore \code{n} e lo passa come argomento alla funzione lambda \code{\textbackslash n -> ...}.
    \item all'interno della lambda viene chiamato \code{return (F n)}; questa funzione prende la foglia \code{F n} e la inserisce in un nuovo costruttore \code{S} che non tocca lo stato, applicando l'azione sul residuo dello stato \code{s'};
    \item lo stato globale aggiornato \code{s'} fluisce così verso l'esterno, pronto per essere ereditato e ulteriormente modificato dalle successive chiamate ricorsive sui rami dell'albero;
\end{enumerate}
\end{gframe}

\subsection{Monad Laws}
Perché una struttura sia una monade valida, le operazioni \code{return} e \code{>>=} devono soddisfare tre leggi fondamentali:

\subsubsection*{Identità Destra}
Applicare \code{return} dopo una computazione non deve alterarne il risultato. "Impacchettare" un valore appena estratto è un'operazione neutra.
\begin{haskellbox}{}
-- notazione standard
p >>= return = p

-- do-notation
do { x <- p; return x } = do { p }
\end{haskellbox}

\subsubsection*{Identità Sinistra}
Se inseriamo un valore in un contesto monadico con \code{return} e lo passiamo a una funzione \code{f} tramite bind, il risultato deve essere identico all'applicazione diretta di \code{f} al valore.
\begin{haskellbox}{}
-- notazione standard
return e >>= f = f e

-- do-notation
do { x <- return e; f x } = do { f e }
\end{haskellbox}

\subsubsection*{Associatività}
L'ordine con cui concateniamo le operazioni non influenza il risultato finale. 

Questa legge permette alla do-notation di comportarsi come una sequenza piatta di comandi.
\begin{haskellbox}{}
-- notazione standard
((p >>= f) >>= g) = p >>= (\x -> (f x >>= g))

-- do-notation
do { y <- do { x <- p; f x }; g y } = do { x <- p; y <- f x; g y }
\end{haskellbox}

\section{Input/Output}
I programmi interattivi sono un problema nel mondo funzionale, perché il valore di una funzione pura dipende solo dai suoi argomenti, mentre input e output sono effettivamente solo dei \textbf{side-effects} (per esempio, se una funzione legge un carattere da tastiera, chiamandola due volte consecutive potremmo ottenere due caratteri diversi anche se non le abbiamo passato alcun argomento diverso).

Per risolvere questo paradosso senza rinunciare alla purezza, Haskell adotta l'idea per cui esiste un argomento implicito che rappresenta lo \textit{stato del mondo}.

Si utilizza quindi una Monade e le espressioni del tipo IO sono dette \textbf{actions}:
\begin{haskellbox}{}
-- I/O vista come una funzione che 'trasforma il mondo'
type IO = World -> World
    
-- in realtà, un'azione IO restituisce anche un valore:
type IO = World -> (a, World)
\end{haskellbox}

La definizione del tipo IO è nascosta al programmatore, in modo da impedirgli di manipolare direttamente \code{World}.

\subsection{Azioni base}
\code{IO Char} rappresenta quindi il tipo delle azioni che ritornano un carattere, e \code{IO ()} quelle che ritornano la tupla vuota (simile a \code{void}), cioè le azioni che sono solo side-effects.

Le primitive di base sono:

\begin{haskellbox}{}
-- per leggere un carattere da tastiera
getChar :: IO Char

-- per scrivere un carattere da tastiera
putChar :: Char -> IO ()

-- per fornire un valore alle azioni 
-- (utile per esempio nei casi base della ricorsione)
return :: a -> IO a
\end{haskellbox}

Se vogliamo sequenzializzare azioni di IO, utilizziamo gli operatori della monade:
\begin{itemize}
    \item se il valore restituito da un'azione non ci interessa (es. \code{putChar}), usiamo \code{>>}
    \item se il valore ci interessa, usiamo \code{>>=} o la \code{do} notation
\end{itemize}

\begin{gframe}{Esempi: scrivere una stringa a video e leggere una riga}
\begin{haskellbox}{}
putStrLn :: String -> IO ()
putStrLn xs = foldr (>>) done (map putChar xs) >> putChar '\n'
\end{haskellbox}
\begin{itemize}
    \item \code{map putChar xs} trasforma una stringa in una lista di azioni \code{IO ()}
    \item \code{foldr (>>)} le concatena in un'unica azione sequenziale, che termina con la stampa del carattere \code{\textbackslash n}
    \item \code{done} corrisponde a \code{return ()}
\end{itemize}   

Sequenziare \code{getChar} è più complicato, perché ci interessa il valore ritornato (il carattere letto).

\begin{haskellbox}{}
getLine = getChar >>= \x ->
            if x == '\n' then return []
            else getLine >>= \xs ->
                return (x:xs)
\end{haskellbox}
\begin{itemize}
    \item \code{getChar} legge un carattere, che viene passato (sequenzialmente da \code{>>=}) alla funzione lambda che controlla se il carattere letto è il newline - se sì, termina la lettura; altrimenti, richiama la funzione ricorsivamente e ritorna il primo carattere letto concatenato al risultato della chiamata
\end{itemize}

\begin{itemize}
    \item \code{getChar} legge un carattere; il suo viene passato tramit \code{>>=} come argomento \code{x} alla prima funzione lambda;
    \item se \code{x} corrisponde al carattere di newline, la lettura termina e viene restituita una lista vuota \code{[]} inserita nel contesto monadico tramite \code{return};
    \item altrimenti, \code{getLine} viene chiamata ricorsivamente per leggere i caratteri rimanenti. Il secondo operatore \code{>>=} estrae la stringa parziale risultante da questa chiamata, legandola alla variabile \code{xs}, e infinre \code{return (x:xs)} ricompone sequenzialmente la stringa completa posizionando il carattere iniziale in testa alla lista;
\end{itemize}
\end{gframe}

\subsection{Strictness dell'I/O}

Anche con le operazioni IO si può usare la \code{do} notation.
\begin{haskellbox}{}
do
  v1 <- a1
  v2 <- a2
  return (f v1 v2)
\end{haskellbox}

La funzione \code{return} è ``a senso unico'': permette di inserire un valore puro in un contesto IO, ma non esiste una funzione simmetrica \code{runIO :: Io a -> a} che permetta di estrarre il valore puro dalla monade.

Se esistesse, si potrebbe scrivere codice come:
\begin{haskellbox}{}
minus :: Int
minus = x - y 
  where
    x = runIO readInt
    y = runIO readInt
\end{haskellbox}

Questo programma sarebbe problematico perché violerebbe il \textbf{Teorema di Church-Rosser} (o di confluenza) che vale per i linguaggi funzionali puri. Il teorema garantisce che l'ordine con cui il compilatore decide di valutare le sotto-espressioni non influenzi mai il risultato finale.

Se dobbiamo calcolare \code{minus x - y}, in una funzione pura il compilatore può calcolare \code{x} e \code{y} nell'ordine in cui preferisce (il risultato sarà identico).

Se invece esistesse \code{runIO}, \code{readInt} chiederebbe all'utente di digitare due numeri. Il compilatore potrebbe decidere di valutare prima \code{x}, e si otterrebbe \code{5 - 3}, oppure prima \code{y}, ottenendo \code{3 - 5}. Lo stesso programma produrrebbe risultati diversi solo a causa dell'ordine di ottimizzazione interna del compilatore. Poiché l'I/O è effectful, l'ordine di lettura è fondamentale.

Per questo, la valutazione nella monade IO è \textbf{strict}.

Per esempio,
\begin{haskellbox}{}
do { undefined; return 3 }
\end{haskellbox}
valuta a \code{undefined} e non a \code{3}, anche se \code{return 3} non avrebbe bisogno dell'eventuale valore calcolato dalla funzione precedente.

\section{State Thread Monad}
Gli algoritmi funzionali puri sulle liste sono eleganti, ma non sempre efficienti. Sebbene molti mantengano la stessa complessità asintotica temporale dei programmi iterativi, la ricorsione introduce un overhead di spazio dovuto ai record di attivazione nello stack. Inoltre, alcune strutture dati sono impossibili da realizzare in modo efficiente senza memoria mutabile: per esempio, le hash table e i dizionari richiedono accesso diretto e modifica in-place degli array.

Per risolvere questo problema, il modulo \code{Control.Monad.ST} fornisce una monade che permette di modificare la memoria in maniera simile ad un linguaggio imperativo.

La struttura di \code{Control.Monad.ST} è analoga a quella della monade State Transformer. Tuttavia, gli elementi della variabile di tipo \code{s} non sono accessibili (come nel caso di \code{World}), se non attraverso le operazioni definite.

\begin{haskellbox}{}
type ST s a = s -> (a, s)
\end{haskellbox}
\begin{itemize}
    \item \code{s} rappresenta quindi una specie di etichetta che identifica un thread di esecuzione (e non si possono trasmettere valori da un thread ad un altro)
\end{itemize}

All'interno di un thread \code{s} si possono creare delle vere e proprie variabili mutabili, chiamate \code{STRef s a} (refeerence a un valore di tipo \code{a} dentro il thread \code{s}).

Le tre primitive esposte da Haskell per lavorare su una variabile \code{STRef} sono:
\begin{itemize}
    \item \code{newSTRef:: a -> ST s (STRef s a)}
        \subitem crea una nuova variabile mutabile, la inizializza con un valore puro e restituisce il suo riferimento
    \item \code{readSTRef :: STRef s a -> ST s a}
        \subitem legge il valore memorizzato nella variabile
    \item \code{writeSTRef :: STRef s a -> a -> ST s ()}
        \subitem sovrascrive il contenuto della variabile (restituisce \code{()} perché serve solo il suo side-effect)
\end{itemize}

\begin{gframe}{Esempio: Fibonacci}
Tramite tupling si può ottenere un Fibonacci lineare, ma che, essendo ricorsivo, occupa $\Theta(n)$ in termini di spazio.

Tramite la monade \code{STRef}, si può ottenere una complessità spaziale di $\Theta(1)$

\begin{haskellbox}{}
-- repeatFor simula un ciclo
repeatFor n = foldr (>>) done . replicate n

fibST :: Int -> ST s Integer
fibST n = do
  a <- newSTRef 0
  b <- newSTRef 1
  repeatFor n (do 
    x <- readSTRef a
    y <- readSTRef b
    writeSTRef a y
    writeSTRef b $! (x + y))
  readSTRef a
\end{haskellbox}

\begin{itemize}
    \item \code{a} e \code{b} vengono allocate ed inizializzate a \code{0} e \code{1}
    \item il blocco interno a \code{repeatFor} viene eseguito $n$ volte; ad ogni iterazione:
        \begin{itemize}
            \item si leggono i valori di \code{a} e \code{b}
            \item si aggiorna il valore di \code{a} con \code{y} e quello di \code{b} con la somma (calcolata in modo strict con \code{\$!}, che forza la valutazione immediata) \code{\$! (x+y)}
            \item alla fine del ciclo, si legge e restituisce il valore accumulato in \code{a}
        \end{itemize}
\end{itemize}
\end{gframe}
Il valore ritornato è però di tipo \code{ST s Integer}. Bisogna quindi estrarre un intero puro.

Serve quindi una funzione \code{runST} analoga a \code{runState} per gli State Transformer.

Questa funzione esiste, ma ha un tipo che non appartiene ai tipi ``standard'' di Haskell (à la Hindley Milner) - contiene infatti un $\forall$ all'interno del tipo invece che all'inizio (apparterrebbe quindi ai tipi del Lambda Calcolo Polimorfo / F2 (vedi \href{https://raw.githubusercontent.com/AglaiaNorza/notes-ig/main/terzo%20anno/linguaggi%20di%20programmazione.pdf}{\underline{appunti del corso ``Linguaggi di programmazione''}})):

\begin{haskellbox}{}
runST :: (forall s. ST s a) -> a
\end{haskellbox}

Un tipo di questo genere si dice \textbf{rank 2} (mentre i tipi di Haskell sono rank 1).

Lo stato deve essere inaccessibile e quindi espressioni come 
\begin{haskellbox}{}
let v = runST (newSTRef True) in runST (readSTRef v)
\end{haskellbox}
non ben tipate. Il type checker non può unificare \code{STRef s a} con un tipo che dipende da \code{s}, quindi non è possibile leggere un valore da un thread e trasportarlo in un altro.

Per estrarre un valore, basta quindi scrivere qualcosa come \code{fib n = runST (fibST n)}.

\subsubsection{Array mutabili}
Un array in Haskell ha tipo \code{STArray s i e}, dove \code{s} è il thread a cui appartiene, \code{i} è il tipo dell'indice, ed \code{e} il tipo dei valori contenuti.

Il tipo \code{i} deve essere istanza della classe \code{i} che contiene tipi enumerabili (con operazioni di \code{succ}) come \code{Int}, \code{Char} ecc.

I metodi fondamentali su \code{STArray} sono:
\begin{itemize}
    \item \code{newListArray (min, max) lista} - inizializza un array mutabile specificando l'intervallo degli indici e inserendo gli elementi di una lista iniziale
    \item \code{readArray xa indice} - legge e restituisce il valore all'indice specificato
    \item{writeArray xa indice valore} - modifica il contenuto di \code{xa} sovrascrivendo il valore all'indice dato con \code{valore} (non ritorna nulla)
    \item{getElems xa} - prende gli elementi dall'array e li copia ``all'indietro'' in una lista haskell
\end{itemize}




\end{document}
